\documentclass[ecta,draft]{econsocart}



% --- DATA FOR PLOTS ---
\begin{filecontents*}{sensitivity_kink.dat}
width mae_kink kink_curvature
64 0.045 12.5
128 0.012 45.2
256 0.003 180.4
512 0.0008 310.1
1024 0.0002 450.5
\end{filecontents*}

\begin{filecontents*}{residual_profile.dat}
wealth residual_early residual_mid residual_converged
0.5 0.1 0.02 0.0001
0.8 0.15 0.03 0.0001
0.95 0.3 0.08 0.0005
1.0 0.8 0.45 0.05
1.05 0.3 0.08 0.0005
1.2 0.15 0.03 0.0001
1.5 0.1 0.02 0.0001
\end{filecontents*}

\begin{filecontents*}{merton_convergence.dat}
epoch c_mean c_std pi_mean pi_std
0.000000 5.443489 0.010000 -0.277846 0.100000
2000.000000 5.270784 0.006703 -0.267854 0.067032
5000.000000 5.357745 0.003679 -0.271296 0.036788
10000.000000 5.430796 0.001353 -0.270625 0.013534
15000.000000 0.030200 0.000700 0.501100 0.004900
19999.000000 0.030200 0.000700 0.501100 0.004900
\end{filecontents*}

\begin{filecontents*}{viscosity_policy.dat}
wealth pi
0.100000 0.513236
0.590000 0.363228
1.080000 0.487027
1.570000 0.534042
2.060000 0.560213
2.550000 0.577754
3.040000 0.589967
3.530000 0.598285
4.020000 0.603704
4.510000 0.607018
5.000000 0.608903
\end{filecontents*}

\begin{filecontents*}{scaling_comparison.dat}
dim time_asg time_nbo error_asg error_nbo
2 0.8 120 1e-9 5e-5
4 15.0 140 1e-7 6e-5
6 420.0 180 5e-6 8e-5
10 12000.0 240 1e-4 1.2e-4
20 NaN 450 NaN 2.5e-4
50 NaN 1100 NaN 5.0e-4
\end{filecontents*}
\begin{filecontents*}{merton_convergence.dat}
epoch c_mean c_std pi_mean pi_std
0.000000 5.443489 0.010000 -0.277846 0.100000
2000.000000 5.270784 0.006703 -0.267854 0.067032
5000.000000 5.357745 0.003679 -0.271296 0.036788
10000.000000 5.430796 0.001353 -0.270625 0.013534
15000.000000 0.030200 0.000700 0.501100 0.004900
19999.000000 0.030200 0.000700 0.501100 0.004900
\end{filecontents*}

\begin{filecontents*}{viscosity_policy.dat}
wealth pi
0.100000 0.513236
0.590000 0.363228
1.080000 0.487027
1.570000 0.534042
2.060000 0.560213
2.550000 0.577754
3.040000 0.589967
3.530000 0.598285
4.020000 0.603704
4.510000 0.607018
5.000000 0.608903
\end{filecontents*}

\begin{filecontents*}{ndu_path_low.dat}
time wealth u pi
0 0.5 4.0 0.10
1 0.52 4.0 0.10
2 0.55 3.9 0.11
3 0.60 3.9 0.11
4 0.68 3.8 0.12
5 0.75 3.8 0.12
6 0.85 3.7 0.13
7 1.0 3.7 0.13
8 1.15 3.6 0.14
9 1.3 3.6 0.14
10 1.5 3.5 0.15
\end{filecontents*}

\begin{filecontents*}{ndu_path_high.dat}
time wealth u pi
0 5.0 4.0 0.25
1 5.5 3.5 0.40
2 6.1 3.0 0.55
3 6.8 2.6 0.65
4 7.7 2.3 0.75
5 8.8 2.1 0.80
6 10.0 2.0 0.85
7 11.5 2.0 0.85
8 13.2 2.0 0.85
9 15.0 2.0 0.85
10 17.0 2.0 0.85
\end{filecontents*}

\begin{filecontents*}{ndu_policy_surface.dat}
1 1 0.05
1 2 0.02
1 3 0.01
1 4 0.01
1 5 0.0
5 1 0.8
5 2 0.6
5 3 0.4
5 4 0.2
5 5 0.1
10 1 1.0
10 2 0.9
10 3 0.7
10 4 0.5
10 5 0.3
\end{filecontents*}

\begin{filecontents*}{ndu_statics.dat}
wealth pi_low_k pi_high_k
1 0.20 0.15
2 0.40 0.25
3 0.60 0.32
4 0.75 0.38
5 0.85 0.42
6 0.90 0.45
7 0.92 0.47
8 0.93 0.48
9 0.94 0.49
10 0.95 0.50
\end{filecontents*}

\begin{filecontents*}{omega_sensitivity.dat}
epoch error_w1 error_w01 error_w10
0 0.1 0.1 0.1
2000 0.01 0.05 0.02
5000 0.001 0.02 0.005
10000 1e-4 0.005 5e-4
15000 1e-5 0.001 1e-4
20000 1e-6 5e-4 5e-5
\end{filecontents*}

\begin{filecontents*}{ti_policy.dat}
wealth c_exp c_soph c_naive
0.5 0.10 0.15 0.20
1.0 0.20 0.28 0.38
1.5 0.30 0.40 0.55
2.0 0.40 0.52 0.70
2.5 0.50 0.64 0.85
3.0 0.60 0.75 1.0
\end{filecontents*}

\begin{filecontents*}{ti_path.dat}
time wealth_exp wealth_soph wealth_naive
0 2.0 2.0 2.0
1 2.2 2.05 1.8
2 2.42 2.10 1.6
3 2.66 2.14 1.4
4 2.9 2.15 1.2
5 3.2 2.16 1.0
\end{filecontents*}

\begin{filecontents*}{cournot_policy_surface.dat}
1 1 10
1 5 8
1 10 6
5 1 15
5 5 12
5 10 9
10 1 20
10 5 16
10 10 13
\end{filecontents*}

\begin{filecontents*}{cournot_comparative_statics.dat}
N K_ss
1 10.0
2 8.0
3 6.5
4 5.5
5 5.0
8 4.0
10 3.5
\end{filecontents*}

\begin{filecontents*}{reference.bib}
@article{barles1991convergence,
  title={Convergence of approximation schemes for fully nonlinear second order equations},
  author={Barles, Guy and Souganidis, Panagiotis E},
  journal={Asymptotic analysis},
  volume={4},
  number={3},
  pages={271--283},
  year={1991},
  publisher={IOS Press}
}

@inproceedings{rahaman2019spectral,
  title={On the spectral bias of neural networks},
  author={Rahaman, Nasim and Baratin, Aristide and Arpit, Devansh and Draxler, Felix and Lin, Min and Hamprecht, Fred and Bengio, Yoshua and Courville, Aaron},
  booktitle={International Conference on Machine Learning},
  pages={5301--5310},
  year={2019},
  organization={PMLR}
}
@article{merton1969lifetime,
  title={Lifetime portfolio selection under uncertainty: The continuous-time case},
  author={Merton, Robert C},
  journal={The Review of Economics and Statistics},
  pages={247--257},
  year={1969},
  publisher={JSTOR}
}
@article{epstein1989,
  title={The structure of preferences, risk, and temporal persistence},
  author={Epstein, Larry G and Zin, Stanley E},
  journal={Econometrica},
  volume={57},
  number={4},
  pages={937--968},
  year={1989}
}
@article{duffie1992stochastic,
  title={Stochastic differential utility},
  author={Duffie, Darrell and Epstein, Larry G},
  journal={Econometrica},
  volume={60},
  number={2},
  pages={353--394},
  year={1992}
}
@book{stokey1989recursive,
  title={Recursive methods in economic dynamics},
  author={Stokey, Nancy L and Lucas, Robert E and Prescott, Edward C},
  year={1989},
  publisher={Harvard university press}
}
@article{raissi2019physics,
  title={Physics-informed neural networks: A deep learning framework for solving forward and inverse problems involving nonlinear partial differential equations},
  author={Raissi, Maziar and Perdikaris, Paris and Karniadakis, George E},
  journal={Journal of Computational Physics},
  volume={378},
  pages={686--707},
  year={2019},
  publisher={Elsevier}
}
@article{han2018solving,
  title={Solving high-dimensional partial differential equations using deep learning},
  author={Han, Jiequn and Jentzen, Arnulf and Weinan, E},
  journal={Proceedings of the National Academy of Sciences},
  volume={115},
  number={34},
  pages={8505--8510},
  year={2018},
  publisher={National Acad Sciences}
}
@article{crandall1983viscosity,
  title={Viscosity solutions of Hamilton-Jacobi equations},
  author={Crandall, Michael G and Lions, Pierre-Louis},
  journal={Transactions of the American Mathematical society},
  volume={277},
  number={1},
  pages={1--42},
  year={1983}
}
@book{fleming2006controlled,
  title={Controlled Markov processes and viscosity solutions},
  author={Fleming, Wendell H and Soner, Halil Mete},
  volume={25},
  year={2006},
  publisher={Springer Science \& Business Media}
}
@book{judd1998numerical,
  title={Numerical methods in economics},
  author={Judd, Kenneth L},
  year={1998},
  publisher={MIT press}
}
@article{brumm2017using,
    title={Using Adaptive Sparse Grids to Solve High-Dimensional Dynamic Models},
    author={Brumm, Johannes and Scheidegger, Simon},
    journal={Econometrica},
    volume={85},
    number={5},
    pages={1575--1616},
    year={2017}
}
@article{hornik1989multilayer,
  title={Multilayer feedforward networks are universal approximators},
  author={Hornik, Kurt and Stinchcombe, Maxwell and White, Halbert},
  journal={Neural networks},
  volume={2},
  number={5},
  pages={359--366},
  year={1989}
}
@article{watkins1992q,
  title={Q-learning},
  author={Watkins, Christopher JCH and Dayan, Peter},
  journal={Machine learning},
  volume={8},
  number={3},
  pages={279--292},
  year={1992}
}
@article{mnih2015human,
    title={Human-level control through deep reinforcement learning},
    author={Mnih, Volodymyr and Kavukcuoglu, Koray and Silver, David and Rusu, Andrei A and Veness, Joel and Bellemare, Marc G and Graves, Alex and Riedmiller, Martin and Fidjeland, Andreas K and Ostrovski, Georg and others},
    journal={Nature},
    volume={518},
    number={7540},
    pages={529--533},
    year={2015},
    publisher={Nature Publishing Group}
}
@book{sutton2018reinforcement,
  title={Reinforcement learning: An introduction},
  author={Sutton, Richard S and Barto, Andrew G},
  year={2018},
  publisher={MIT press}
}
@article{kingma2014adam,
    title={Adam: A method for stochastic optimization},
    author={Kingma, Diederik P and Ba, Jimmy},
    journal={arXiv preprint arXiv:1412.6980},
    year={2014}
}
@article{hornik1990universal,
    title={Universal approximation of an unknown mapping and its derivatives using multilayer feedforward networks},
    author={Hornik, Kurt and Stinchcombe, Maxwell and White, Halbert},
    journal={Neural Networks},
    volume={3},
    number={5},
    pages={551--560},
    year={1990}
}
@book{borkar2008stochastic,
    title={Stochastic approximation: a dynamical systems viewpoint},
    author={Borkar, Vivek S},
    year={2008},
    publisher={Cambridge University Press}
}
@article{mcfadden1989method,
    title={A method of simulated moments for estimation of discrete response models without numerical integration},
    author={McFadden, Daniel},
    journal={Econometrica},
    volume={57},
    number={5},
    pages={995--1026},
    year={1989}
}
@article{duffie1993simulated,
    title={Simulated moments estimation of Markov models of asset prices},
    author={Duffie, Darrell and Singleton, Kenneth J},
    journal={Econometrica},
    volume={61},
    number={4},
    pages={929--952},
    year={1993}
}
@article{laibson1997golden,
    title={Golden eggs and hyperbolic discounting},
    author={Laibson, David},
    journal={The Quarterly Journal of Economics},
    volume={112},
    number={2},
    pages={443--478},
    year={1997}
}
@article{harris2003stochastic,
    title={Stochastic equilibria in a time-inconsistent optimal problem},
    author={Harris, Christopher and Laibson, David},
    year={2003},
    journal={Unpublished manuscript, Cambridge University}
}
@article{ericson1995markov,
    title={Markov-perfect industry dynamics: A framework for empirical work},
    author={Ericson, Richard and Pakes, Ariel},
    journal={The review of economic studies},
    volume={62},
    number={1},
    pages={53--82},
    year={1995}
}
@article{mishra2022error,
    title={Error estimates for physics informed neural networks for linear elliptic PDEs},
    author={Mishra, Siddhartha and Molinaro, Roberto},
    journal={IMA Journal of Numerical Analysis},
    volume={42},
    number={4},
    pages={2972--3013},
    year={2022}
}
@article{karniadakis2021physics,
    title={Physics-informed machine learning},
    author={Karniadakis, George Em and Kevrekidis, Ioannis G and Lu, Lu and Perdikaris, Paris and Wang, Sifan and Yang, Liu},
    journal={Nature Reviews Physics},
    volume={3},
    number={6},
    pages={422--440},
    year={2021}
}
@article{qi2025,
  title={Neural Foundations of Endogenous Risk Aversion},
  author={Qi, Qian},
  journal={Unpublished Manuscript, Peking University},
  year={2025}
}
@article{qi2025universalappr,
    title={Universal Approximation Theorems for Differentiable Deep Reinforcement Learning},
    author={Qi, Qian},
    journal={Unpublished Manuscript, Peking University},
    year={2025}
}
@inproceedings{qian2025icml,
    title={Deep Learning is a Good Basis Function},
    author={Qian, Qi},
    booktitle={International Conference on Machine Learning},
    year={2025},
    organization={PMLR}
}
\end{filecontents*}

% --- PREAMBLE AND PACKAGES ---
\RequirePackage[colorlinks,citecolor=blue,linkcolor=blue,urlcolor=blue,pagebackref]{hyperref}
\usepackage{setspace,graphicx,epstopdf,amsmath,amsfonts,amssymb,amsthm}
\usepackage{marginnote,datetime,enumitem,subfigure,rotating,fancyvrb, booktabs}
\usepackage{float}
\usdate
\usepackage{lscape}
\usepackage{adjustbox}
\usepackage{comment}
\usepackage{multirow}
\usepackage{pdflscape}
\usepackage{tikz}
\usepackage{caption}
\usepackage{xcolor}
\usepackage{algorithm}
\usepackage{algorithmicx}
\usepackage{algpseudocode}
\usepackage{longtable}
\usepackage{threeparttable}
\usepackage{tabularx}
\usepackage{cleveref}
\usepackage{ragged2e}
\usepackage{mathpazo,newpxtext}
\usepackage[T1]{fontenc}
\usepackage{bm} % For bold math symbols
\usepackage{etoolbox}
\sloppy

\usepackage{pgfplots}
\pgfplotsset{compat=1.18}
\usepgfplotslibrary{groupplots}
\usepgfplotslibrary{fillbetween}
\usepgfplotslibrary{statistics}
\pgfplotsset{
    every axis/.append style={
        label style={font=\small},
        tick label style={font=\small},
        title style={font=\bfseries}
    },
    legend style={font=\small}
}


% --- LOCAL DEFINITIONS ---
\startlocaldefs

\theoremstyle{plain}
\newtheorem{theorem}{Theorem}
\newtheorem{lemma}[theorem]{Lemma}
\newtheorem{proposition}{Proposition}
\newtheorem{assumption}{Assumption}

\theoremstyle{remark}
\newtheorem{definition}{Definition}
\newtheorem{example}{Example}
\newtheorem{remark}{Remark}

\newcommand{\E}{\mathbb{E}}
\newcommand{\R}{\mathbb{R}}
\DeclareMathOperator*{\argmax}{arg\,max}
\DeclareMathOperator*{\argmin}{arg\,min}
\newcommand{\hamiltonian}{\mathcal{H}}
\newcommand{\statevec}{\bm{S}}
\newcommand{\controlvec}{\bm{\alpha}}
\newcommand{\driftvec}{\bm{\mu}}
\newcommand{\diffvec}{\bm{\sigma}}
\newcommand{\paramvec}{\bm{\theta}}

\endlocaldefs

% --- DOCUMENT START ---
\begin{document}

\begin{frontmatter}

\title{Neural Bellman Operators}
\runtitle{Neural Bellman Operators}

\begin{aug}
\author[id=au1,addressref={add1}]
{\fnms{Qian}~\snm{QI}\ead[label=e1]{qiqian@pku.edu.cn}}

\address[id=add1]{%
\orgdiv{Peking University},
\orgname{No.5 Yiheyuan Road Haidian District, Beijing, P.R.China 100871}}

\end{aug}

\support{This paper serves as a numerical method companion to the initial contribution \cite{qi2025}. I extend my gratitude to Darrell Duffie for his insightful and very detailed discussions. I also thank the editor and an anonymous referee for their exceptionally constructive feedback, which has significantly improved the rigor and clarity of this manuscript.}

\begin{abstract}
High-dimensional dynamic stochastic models, central to modern economic theory, remain largely intractable. This paper introduces the Neural Bellman Operators (NBO) method, a scalable algorithm for solving such problems. The method employs an actor-critic framework where the critic is trained to minimize the residual of the Hamilton-Jacobi-Bellman (HJB) equation. This HJB-based objective provides a stable, zero-centered learning target that robustly handles the recursive utility structures that destabilize standard reinforcement learning algorithms. Furthermore, the architecture is uniquely suited to find the non-differentiable viscosity solutions that characterize models with inequality constraints. We validate NBO by recovering known analytical and viscosity solutions and apply it to a high-dimensional model with endogenous risk aversion, uncovering novel dynamics. While the underlying algorithm is formally guaranteed to find a local equilibrium, we provide robust empirical evidence of convergence to correct, globally optimal solutions. NBO provides a flexible and powerful tool for bringing complex dynamic theories to quantitative analysis.
\end{abstract}

\begin{keyword}
\kwd{Deep Learning}
\kwd{Stochastic Optimal Control}
\kwd{Recursive Utility}
\kwd{HJB Equation}
\kwd{Viscosity Solution}
\kwd{High-Dimensional Models}
\end{keyword}

\end{frontmatter}

\section{Introduction}

This paper introduces a general-purpose computational methodology, the Neural Bellman Operators (NBO) method, for solving high-dimensional continuous-time stochastic optimal control problems in economics. Our work is motivated by the challenge of bringing modern economic theories to quantitative analysis. Many frontier models in macroeconomics and finance feature high-dimensional state spaces that are computationally infeasible with traditional grid-based methods. This \textit{curse of dimensionality} is a long-standing barrier not only for models of endogenous preferences with joint optimization problems (see \Cref{eq:gsc}, also \cite{qi2025}), but also for large-scale heterogeneous agent models, optimal Ramsey policy with complex state dynamics, and dynamic games.


A deeper structural challenge, central to modern economic theory, is the recursive nature of the agent's objective. In models with recursive utility (e.g., \cite{epstein1989,duffie1992stochastic}) or endogenous preferences (\cite{qi2025}), the value function $V_s$ becomes a direct argument in the flow utility function, $f(s, \statevec_s, \theta_s, a_s, V_s)$. This feature poses a significant problem for many standard reinforcement learning (RL, see 
\cite{watkins1992q,mnih2015human,sutton2018reinforcement}) algorithms. Such methods typically rely on an additively separable reward structure. When flow utility depends on the value function, the Bellman target used to learn the value function, for instance, the discrete-time expression $f(s_t, a_t, V(s_t))\Delta t + V(s_{t+1})$, becomes a function of the object being learned at the current state, $V(s_t)$. This creates a difficult \textit{moving target} problem that can destabilize or slow down the learning process. An effective solver must therefore not only scale to high dimensions but also be structurally equipped to handle this recursive dependency and find the non-differentiable \textit{viscosity solutions} that arise from economic constraints.

This paper bridges this gap by developing the \textbf{Neural Bellman Operators (NBO)}, a method uniquely structured to find these solutions. At its core, NBO extends the logic of recursive methods (see \cite{stokey1989recursive}) by parameterizing the value and policy functions with neural networks and training them to find a fixed point of the continuous-time Bellman operator. It achieves this through a novel actor-critic mechanism that finds a Nash equilibrium between a constrained \textit{actor} network that learns the optimal policy within the feasible set, automatically identifying any \textit{kinks} at the boundaries, and a \textit{critic} network that learns the corresponding non-differentiable value function by enforcing the HJB equation. The term \textbf{Bellman Operator} refers to the right-hand side of the HJB equation, which generates the Bellman dynamics. Our method works by ensuring the neural network approximations satisfy the conditions defined by this operator. Our approach builds upon recent breakthroughs in solving high-dimensional PDEs, adapting ideas from Physics-Informed Neural Networks (PINNs, see \cite{raissi2019physics}) and deep BSDE methods (e.g., \cite{han2018solving}) into a novel framework tailored for economic control problems. The core idea is to directly approximate the agent's optimal policies and the associated value function using deep neural networks and optimize them to satisfy the underlying optimality conditions from control theory.

This method offers several key advantages. First and foremost, it directly handles recursive utility. By integrating the Hamilton-Jacobi-Bellman (HJB) residual directly into an actor-critic loop, our method provides a stable, zero-centered learning target for the value function (the \textit{critic}) even when it is an argument in the flow utility. This robustly addresses the recursive challenge that complicates many standard RL algorithms. Second, its computational complexity scales polynomially, not exponentially, with state-space dimensionality, making it scalable to problems previously considered intractable. Third, the framework is general. By leveraging the recent universal approximation theoretical guarantees of deep reinforcement learning on viscosity solutions (see \cite{qi2025universalappr,qian2025icml}), we provide strong empirical evidence that NBO robustly approximates the non-differentiable \textit{viscosity solutions} that arise in models with inequality constraints, significantly broadening its practical applicability. Finally, a single algorithmic core can solve a wide variety of economic problems, including models with time-inconsistent agents and dynamic games, by simply specifying the relevant state dynamics and objectives. The practical utility of the method rests on compelling empirical evidence that the algorithm, despite navigating a non-convex loss landscape, consistently finds high-quality solutions that match known analytical results and produce novel, economically intuitive outcomes in more complex settings.

\subsection{Related Literature}
Our work sits at the intersection of three distinct fields: computational economics, machine learning for solving differential equations, and reinforcement learning. The central challenge addressed by this paper, the curse of dimensionality in dynamic programming, is a foundational problem in computational economics. Since the formalization of recursive methods (see \cite{stokey1989recursive}), the standard approach for solving such problems has been grid-based value or policy function iteration. While powerful, these methods require discretizing the state space, and their computational cost grows exponentially with the number of state variables, rendering them impractical for the high-dimensional problems that motivate our work. A vast literature on numerical methods has sought to overcome this barrier using approximation techniques such as projection methods and parameterized expectations, which approximate the solution with a linear combination of basis functions (e.g., \cite{judd1998numerical}). Our approach follows this approximation philosophy but leverages the power of deep neural networks as universal function approximators (e.g., \cite{hornik1989multilayer,qian2025icml}), which do not require the user to specify basis functions and are more naturally suited to high-dimensional spaces.

\begin{sloppypar}
The technical core of our method for learning the value function is drawn from recent advances in using neural networks to find solutions to partial differential equations. The key idea of Physics-Informed Neural Networks (PINNs) is to include the governing PDE's residual as a penalty term in the loss function, thereby forcing the network's output to conform to the underlying physical or mathematical laws (see \cite{raissi2019physics}). Our HJB loss ($L_{\text{HJB}}$) is a direct application of this principle, where the HJB equation represents the law that the value function must obey. This approach is closely related to methods that use deep learning to solve Forward-Backward Stochastic Differential Equations (FBSDEs), such as the Deep BSDE method of \cite{han2018solving}. However, as we detail in Section \ref{sec:relation_to_bsde}, the NBO method is fundamentally distinct because it is designed to solve a control problem directly, rather than a pre-specified FBSDE system. NBO automates the discovery of the optimal feedback control rule, bypassing the often-intractable analytical pre-computation of the FBSDE system required by related methods.
\end{sloppypar}

Finally, the overall architecture of our algorithm is an actor-critic method, a cornerstone of modern reinforcement learning (e.g., \cite{sutton2018reinforcement}). In this framework, an \textit{actor} (our policy networks $\Theta_\psi, \mathcal{A}_\phi$) learns how to act, while a critic (our value network $V_\xi$) learns to evaluate the actor's policy. Our contribution is to design a \textit{critic} specifically for continuous-time economic models. Standard reinforcement learning algorithms, which are often developed for discrete-time problems with additively separable rewards, can become unstable when applied to models with recursive utility, where the flow utility $f$ depends on the value function $V_t$. This introduces a problematic self-reference in the learning target. The NBO method's critic, by contrast, is trained to minimize the HJB residual, a stable, zero-centered objective that does not suffer from this \textbf{moving target} problem.


\paragraph*{High-Performance Computing and Adaptive Grids.} It is essential to contrast our approach with the other dominant paradigm for solving high-dimensional dynamic models: the use of massively parallel High-Performance Computing (HPC) combined with sophisticated projection methods (see \cite{brumm2017using}). This approach attacks the curse of dimensionality not by abandoning grids, but by making them more intelligent. Using techniques like adaptive sparse grids, it represents the solution with a basis of global functions (e.g., polynomials) on a cleverly chosen, small set of grid points. The problem is then transformed into a massive system of algebraic equations solved on supercomputers.

The NBO method represents a fundamentally different philosophy. It is a \textit{mesh-free} method that learns the solution by stochastically sampling from the state space, rather than relying on a deterministic grid. Instead of using pre-specified basis functions, NBO employs neural networks to learn the necessary functional form directly from the problem's HJB constraint. These paradigms should be viewed as complementary. HPC methods may be more efficient for problems with known smoothness properties where good basis functions can be chosen, whereas NBO is particularly advantageous for problems where the solution's functional form is unknown or highly non-smooth. Consequently, while the HPC approach is designed for massively parallel infrastructure, NBO is designed to be effective on a single machine with a standard GPU, offering greater accessibility.


The primary contribution of this paper is the synthesis of these three streams of literature into a single, cohesive framework. We employ a PINN-style loss function as a novel and stable learning objective for the critic in an actor-critic algorithm, and we deploy this architecture to solve the high-dimensional dynamic programming problems that have long been a frontier in quantitative economics.

This paper is organized as follows. \Cref{sec:problem_formulation} formalizes the general constrained control problem and its HJB representation. \Cref{sec:nbo_method} presents the NBO algorithm with its theoretically grounded loss function. \Cref{sec:theoretical_framework} discusses the theoretical framework and its scope, with a focus on viscosity solutions. \Cref{sec:validation} provides essential numerical validation, including on a problem with a known viscosity solution. \Cref{sec:recursive_test} validates the method on a model with fully recursive utility. \Cref{sec:application_ndu} presents novel economic insights from an NDU model of portfolio choice. \Cref{app:extensions} demonstrates the framework's flexibility by outlining solutions for models with time-inconsistent agents and dynamic games. Finally, \Cref{sec:conclusion} concludes. Appendices provide implementation details, hyperparameter analysis, formal proofs, and discussion of extensions.

\section{The General Stochastic Control Problem}
\label{sec:problem_formulation}

We define a generalized stochastic control problem that encompasses NDU models (see \cite{qi2025}) as a special case, driven by a general martingale process.

\begin{definition}[The General Stochastic Control Problem]\label{eq:gsc}
Let $(\Omega, \mathcal{F}, \{\mathcal{F}_t\}_{t \ge 0}, \mathbb{P})$ be a filtered probability space. An agent's decision problem is characterized by a state space $(u_t, X_t) \in \R^q \times \R^k$ and a control space $(\theta_t, a_t) \in \Theta \times \mathcal{A}$. Let $\statevec_t = (u_t, X_t)$ be the full state vector. The agent chooses admissible, progressively measurable control paths $\{\theta_s, a_s\}_{s=t}^T$ to maximize the objective functional:
\begin{equation}
\label{eq:objective}
V(t, \statevec_t) = \sup_{\{\theta_s, a_s\}} \E_t \left[ \int_t^T \left( f(s, \statevec_s, \theta_s, a_s, V(s, \statevec_s)) - \lambda \mathcal{L}(u_s, \theta_s) \right) ds + G(\statevec_T) \right]
\end{equation}
subject to the forward Stochastic Differential Equations (SDEs):
\begin{align}
du_t &= h(t, \statevec_t, \theta_t)dt + \sigma^h(t, \statevec_t, \theta_t) dM_t^h, & u_0 = u_{init} \label{eq:state_u} \\
dX_t &= \mu_X(t, \statevec_t, a_t)dt + \sigma_X(t, \statevec_t, a_t) dM_t^X, & X_0 = x_{init} \label{eq:state_x}
\end{align}
where $V_s \equiv V(s, \statevec_s)$ is the value function. The term $f(\cdot, V_s)$ allows for recursive utility. $\mathcal{L}$ is a cost function, $G$ is a terminal utility function, and $M_t^h, M_t^X$ are components of a vector-valued, continuous $\mathcal{F}_t$-martingale $\bm{M}_t$, which may have correlated components.
\end{definition}

\subsection{Hamilton-Jacobi-Bellman Representation}
This stochastic optimal control problem is characterized by a Hamilton-Jacobi-Bellman (HJB) partial differential equation. To state this formally, we first define the pre-supremum Hamiltonian operator for a given policy $(\theta, a)$, which generates the dynamics of the value function under that policy:
\begin{equation}
\mathcal{H}^{\theta,a}(t, \statevec, V, \nabla V, \nabla^2 V) := f(t,\statevec,\theta,a,V) - \lambda \mathcal{L}(u, \theta) + \mathcal{D}^{\theta,a}V(t,\statevec).
\end{equation}
The term $\mathcal{D}^{\theta,a}V$ is the Itô differential operator applied to $V(t, \statevec_t)$ under controls $(\theta,a)$. Let $\statevec_t = (u_t, X_t)$, and let the combined drift and diffusion of the state vector be $\driftvec_{\statevec}(t, \statevec_t, \theta, a)$ and $\diffvec_{\statevec}(t, \statevec_t, \theta, a)$ respectively. Let the quadratic variation process of the martingale $\bm{M}_t$ be given by $d\langle \bm{M} \rangle_t = \bm{Q}_t dt$, where $\bm{Q}_t$ is a predictable matrix-valued process. The operator expands as:
\begin{equation}
\mathcal{D}^{\theta,a}V = (\nabla_{\statevec} V)^\top \driftvec_{\statevec} + \frac{1}{2}\text{Tr}\left(\diffvec_{\statevec} \bm{Q}_t \diffvec_{\statevec}^\top \nabla_{\statevec\statevec}^2 V\right),
\end{equation}
where $\nabla_{\statevec} V$ is the gradient and $\nabla_{\statevec\statevec}^2 V$ is the Hessian matrix of $V$. When $\bm{M}_t$ is a standard vector Brownian motion, $\bm{Q}_t$ is simply the identity matrix.

Assuming sufficient smoothness, the value function $V(t,\statevec_t)$ must then satisfy the HJB equation, which states that the time decay of the value must equal the value generated by the optimal policy:
\begin{align}
\label{eq:hjb}
-\partial_t V(t,\statevec_t) &= \sup_{\theta, a} \mathcal{H}^{\theta,a}(t, \statevec_t, V, \nabla_{\statevec} V, \nabla_{\statevec\statevec}^2 V) \nonumber\\&\equiv \hamiltonian(t, \statevec_t, V, \nabla_{\statevec} V, \nabla_{\statevec\statevec}^2 V),
\end{align}
with the terminal condition $V(T, \statevec_T) = G(\statevec_T)$. The term \textbf{Neural Bellman Operators} refers to our method of using neural networks to find a fixed point where the learned value function and policy satisfy this optimality condition. The accurate computation of the Hessian via automatic differentiation is a core feature of our method. Solving this high-dimensional PDE directly is the source of the curse of dimensionality.

\section{The Neural Bellman Operators Method (NBO)}
\label{sec:nbo_method}

NBO approximates the optimal policies and the value function with neural networks. The key innovation is a composite loss function that forces the networks to find a fixed point of the Bellman operator by satisfying the HJB equation.

\subsection{Network Parameterization and Constraint Handling}
We employ three distinct deep neural networks, each mapping the current state $(t, \statevec_t)$ to a specific component of the solution. An internal policy network, $\Theta_{\psi}(t, \statevec_t)$ with weights $\psi$, outputs the internal controls $\theta_t$. An external policy network, $\mathcal{A}_{\phi}(t, \statevec_t)$ with weights $\phi$, outputs the external actions $a_t$. Finally, a value function network, $V_{\xi}(t, \statevec_t)$ with weights $\xi$, outputs an estimate of the value function $V_t$. Collectively, these networks represent an approximation to the full solution of the control problem. The optimization problem is to find weights $(\psi^*, \phi^*, \xi^*)$ that solve the problem.

A crucial aspect of solving economic problems is respecting constraints on control variables (e.g., non-negativity of investment or consumption, borrowing limits). The NBO framework handles these constraints directly within the actor's optimization. This feature is not an incidental detail; it is the fundamental mechanism that enables the method to find non-differentiable viscosity solutions, as discussed in Section \ref{sec:theoretical_framework}. The method for enforcing constraints depends on their structure:
\begin{enumerate}
    \item \textbf{Box Constraints:} For simple constraints of the form $a_t \in [a_{min}, a_{max}]$, applying a bounded activation function to the network's output layer is efficient and effective. For example, using the output $a_t = a_{min} + (a_{max} - a_{min}) \cdot \text{sigmoid}(\text{NN}_{\text{output}})$ ensures the constraint is met by construction. All experiments in this paper fall into this category.
    \item \textbf{General Constraints:} For more complex, possibly state-dependent, constraints of the form $g(\statevec_t, a_t) \le 0$, a penalty-based approach is required. The actor's performance loss $L_{\text{PERF}}$ can be augmented with a penalty term, for example: $L_{\text{PERF}}^{\text{constr}} = L_{\text{PERF}} + \lambda_g \E[\max(0, g(\statevec_t, a_t))^2]$, where $\lambda_g$ is a large penalty coefficient. This transforms the constrained problem into an unconstrained one that the NBO algorithm can solve, extending its applicability to a wide range of economic models.
\end{enumerate}

\subsection{The Algorithm and Loss Function}
The algorithm iteratively improves the networks. The crucial element is the loss function, which consists of two parts, embodying the principles of a PINN-constrained actor-critic algorithm.

\begin{definition}[The NBO Loss Function]
The total loss to be minimized with respect to parameters $(\psi, \phi, \xi)$ is a weighted sum of two components:
\begin{equation}
\label{eq:loss_total}
\text{Loss}(\psi, \phi, \xi) = L_{\text{PERF}}(\psi, \phi, \xi) + \omega L_{\text{HJB}}(\psi, \phi, \xi),
\end{equation}
where $\omega > 0$ is a hyperparameter weighting the two objectives.

The first component, the performance loss $L_{\text{PERF}}$ (the \textit{actor} loss), drives the policy networks to maximize the rate of value accumulation at each state-time point. This rate is precisely the Hamiltonian, which combines the flow utility with the expected change in continuation value (as measured by the critic). The actor's objective is therefore to maximize the expected Hamiltonian, which is equivalent to minimizing its negative:
\begin{equation}
\label{eq:loss_perf}
    L_{\text{PERF}} = - \E \left[ \hamiltonian^{\psi,\phi}(t, \statevec_t, V_\xi, \nabla_{\statevec} V_\xi, \nabla_{\statevec\statevec}^2 V_\xi) \right],
\end{equation}
where $\hamiltonian^{\psi,\phi}$ is the pre-supremum Hamiltonian evaluated at the current policies $\theta_t = \Theta_\psi(t,\statevec_t)$ and $a_t = \mathcal{A}_\phi(t,\statevec_t)$, and using the critic's current value function $V_\xi$ and its derivatives. This formulation provides a direct, local, and low-variance gradient signal for policy improvement.

The second component, the HJB consistency loss $L_{\text{HJB}}$ (the \textit{critic} loss), is the critic's objective. It penalizes deviations from the Bellman equation, forcing the value network $V_\xi$ to be a consistent valuation of the policies currently generated by the actor networks $\Theta_\psi$ and $\mathcal{A}_\phi$:
\begin{equation}
\label{eq:loss_hjb}
    L_{\text{HJB}} = \E \left[ \left( -\partial_t V_\xi(t,\statevec_t) - \mathcal{H}^{\psi,\phi}(t, \statevec_t, V_\xi, \nabla_{\statevec} V_\xi, \nabla_{\statevec\statevec}^2 V_\xi) \right)^2 \right]. 
\end{equation}
The gradients $\partial_t V_\xi, \nabla_{\statevec} V_\xi$, and Hessian $\nabla_{\statevec\statevec}^2 V_\xi$ are computed via automatic differentiation of the network $V_\xi$. The expectation is taken over a batch of randomly sampled state-time points $(t, \statevec_t)$ from the problem domain.
\end{definition}

The composite loss function in Eq. \eqref{eq:loss_total} establishes a game between the policy networks (the \textit{actor}) and the value network (the \textit{critic}). The actor-critic interpretation is central: $L_{\text{HJB}}$ is the critic's learning rule. For any policy supplied by the actor, the critic's objective is to minimize this loss, forcing it to become a faithful reporter of the true value of that policy. The actor, in turn, takes the critic's valuation (i.e., $V_\xi$ and its derivatives) as given and seeks to find the policy that maximizes the Hamiltonian by minimizing $L_{\text{PERF}}$. This design replaces the explicit $\sup$ operator in the HJB equation with a smooth gradient-based maximization, allowing the policies and the value function to co-evolve toward the optimal solution.

\begin{remark}[The Actor's Objective and the Policy Gradient]
\label{rem:actor_objective}
The formulation of the actor's loss $L_{\text{PERF}}$ as the negative expected Hamiltonian is the continuous-time analogue of the policy gradient theorem in actor-critic methods (\cite{sutton2018reinforcement}). The agent's objective is to find a policy $\controlvec^{\bm{\pi}}$ with parameters $\bm{\pi}=(\psi, \phi)$ that maximizes the lifetime objective, $J(\bm{\pi}) = V^{\bm{\pi}}(0, \statevec_0)$. The Policy Gradient Theorem establishes that the gradient of this objective with respect to the policy parameters can be expressed in terms of the Hamiltonian. For a given state, the change in value with respect to a change in policy is driven by how that policy change affects the Hamiltonian. The gradient ascent update for the policy parameters is thus proportional to the gradient of the Hamiltonian with respect to those parameters:
\begin{equation}
\bm{\pi}_{k+1} \leftarrow \bm{\pi}_k + \eta \nabla_{\bm{\pi}} \mathcal{H}^{\bm{\pi}}(t_k, \statevec_k, V^{\bm{\pi}}, \dots).
\end{equation}
Minimizing our defined $L_{\text{PERF}}$ achieves precisely this gradient ascent step. The critic provides an estimate of the value function and its derivatives, $\hat{V} \approx V^{\bm{\pi}}$, and the actor performs a gradient step to maximize the Hamiltonian $\mathcal{H}^{\bm{\pi}}(t, \statevec, \hat{V}, \dots)$. This ensures that the actor is always taking steps to improve its policy based on the critic's current best estimate of the consequences. This formulation has significantly lower variance than using the gradient of the integrated lifetime utility, leading to more stable learning.
\end{remark}

A subtle but crucial mechanism arises in recursive models where the flow utility $f$ is a function of the value $V$. In this case, the total gradient for the critic's parameters $\xi$ contains two terms: $\nabla_{\xi}L_{\text{PERF}}$ and $\omega\nabla_{\xi}L_{\text{HJB}}$. While the HJB loss is the primary driver for learning the value function's dynamic properties (i.e., its shape, governed by its derivatives), the gradient from the performance loss plays a vital role in anchoring the function's absolute \textit{level}. To see this, note that the HJB loss primarily constrains the relationship between the derivatives of $V_\xi$. The performance loss, in contrast, contains the flow utility $f(\cdot, V_\xi)$. Its gradient with respect to the critic's parameters,
\begin{equation}
\nabla_{\xi} L_{\text{PERF}} = - \E \left[ \frac{\partial f}{\partial V}(t,\statevec_t,a_t,V_\xi) \cdot \nabla_{\xi} V_\xi(t, \statevec_t) \right],
\end{equation}
provides a direct, global signal to shift the entire value function level up or down to maximize total lifetime utility.


The stability of the NBO method in the presence of recursive utility is a direct and crucial consequence of this architecture. To fully appreciate this advantage, it is instructive to contrast our HJB-based critic with a more conventional approach from reinforcement learning based on a naive temporal-difference (TD) error. Consider a discrete-time approximation of the Bellman equation for a recursive problem. The value of a policy at time $t$ is the expected sum of the flow utility from $t$ to $t+\Delta t$ and the continuation value at $t+\Delta t$. For a recursive utility function $f(s, a, V)$, this is:
\begin{equation}
    V(s_t) \approx \E_t \left[ f(s_t, a_t, V(s_t))\Delta t + V(s_{t+1}) \right].
\end{equation}
A standard TD algorithm would train the value network $V_\xi$ by minimizing the squared \textbf{TD error}, which is the difference between a target $Y_t$ and the current value estimate $V_\xi(s_t)$. The target, constructed from the right-hand side of the equation above, would be:
\begin{equation}
\label{eq:td_target}
    Y_t = f(s_t, a_t, V_\xi(s_t))\Delta t + V_\xi(s_{t+1}).
\end{equation}
The resulting TD error to be minimized is $\delta_t = Y_t - V_\xi(s_t)$. The critical issue is now apparent: when we substitute the target, the error term becomes:
\begin{equation}
\label{eq:td_error}
    \delta_t = \left[ f(s_t, a_t, V_\xi(s_t))\Delta t + V_\xi(s_{t+1}) \right] - V_\xi(s_t).
\end{equation}
The network's output at state $s_t$, $V_\xi(s_t)$, appears on both sides of the subtraction. The learning target for $V_\xi(s_t)$ is an explicit function of itself. This creates a problematic self-referential dependency, known as a \textbf{moving target} problem. Each time the network's weights $\xi$ are updated to reduce the error, the target itself shifts, which can lead to unstable oscillations or outright divergence during training.

The NBO method's critic, by contrast, is trained to minimize the HJB consistency loss $L_{\text{HJB}}$, which sidesteps this problem entirely. The learning objective is to drive the HJB residual to zero across the state-time domain:
\begin{equation}
    \text{Residual}(t,\statevec_t) = -\partial_t V_\xi(t,\statevec_t) - \mathcal{H}^{\psi,\phi}(t, \statevec_t, V_\xi, \nabla_{\statevec} V_\xi, \nabla_{\statevec\statevec}^2 V_\xi) \longrightarrow 0
\end{equation}
Even though the Hamiltonian $\mathcal{H}^{\psi,\phi}$ is a function of $V_\xi$ in a recursive model, the \textit{target} for the entire residual expression is identically zero. This provides a fixed, stationary learning objective that does not depend on the level of the value function itself. The optimization goal is to make the learned dynamics of the value function (its partial derivatives) consistent with the economics of the problem, rather than chasing a moving value target. This architectural choice is the key to the algorithm's stability in recursive settings, a claim we empirically validate in our application to Epstein-Zin preferences in Section \ref{sec:recursive_test}.

It is important to highlight the interdependence inherent in this formulation. The actor's objective, $L_{\text{PERF}}$, is evaluated using the critic's current value estimate, $V_\xi$. This means the policy networks are optimizing against a target provided by a critic that is itself learning. The stability of the NBO framework hinges on the nature of the critic's loss, $L_{\text{HJB}}$. Because the HJB residual provides a fixed, zero-centered target that is independent of the absolute level of the value function, it offers a stable learning signal for the critic. This, in turn, provides a more reliable evaluation for the actor, mitigating the instabilities that can arise from such interdependent learning and ensuring robust joint convergence. As detailed in the proof of Proposition \ref{prop:optimality} (Appendix \ref{app:proofs}), the joint minimization seeks a Nash equilibrium of this game between the actor and critic, which corresponds to the true optimal solution of the original control problem.

\subsection{The Actor-Critic Balance and Hyperparameter \texorpdfstring{$\omega$}{omega}}
\label{sec:omega}
The hyperparameter $\omega$ in the loss function \eqref{eq:loss_total} is critical as it balances policy optimization (the actor's goal) against Bellman consistency (the critic's goal). In the game-theoretic framing, it can be interpreted as setting the relative learning speed or \textit{power} in the game between the actor and critic. An appropriate balance is necessary for stable convergence. If $\omega$ is too small, the algorithm may prioritize finding policies that perform well under an inaccurate value function that does not satisfy the HJB equation. If $\omega$ is too large, the algorithm can become overly focused on minimizing the HJB residual for potentially suboptimal policies, thereby stifling exploration. From a theoretical perspective (see Appendix \ref{app:convergence}), the choice of $\omega$ can be interpreted as setting the relative learning speed between the critic and the actor, with a larger $\omega$ promoting the "fast" critic convergence that is crucial for stable actor-critic learning. A value of $\omega=1.0$ serves as a robust default because it tends to normalize the gradients from the performance loss and the HJB loss to similar orders of magnitude in canonical problems. This promotes a balanced learning rate between the actor and the critic. For this reason, and to demonstrate the method's robustness, we use $\omega=1.0$ for all experiments in this paper. A detailed sensitivity analysis validating this choice is provided in Appendix \ref{app:hyperparams}.

\subsection{Relation to Existing Methods: NBO vs. Deep BSDE}
\label{sec:relation_to_bsde}

The NBO method represents a synthesis of ideas from Physics-Informed Neural Networks (PINNs) and actor-critic reinforcement learning. While it shares the use of deep learning to solve high-dimensional PDEs with the Deep BSDE method of \cite{han2018solving}, the two approaches differ fundamentally in their mathematical formulation, algorithmic complexity, and domain of economic applicability.

The Deep BSDE method is a solver for \textit{Forward-Backward Stochastic Differential Equations (FBSDEs)}. In contrast, NBO is a direct solver for the \textit{primal stochastic control problem} via the HJB residual. This distinction dictates a clear trade-off between computational complexity and structural flexibility.

\paragraph{Semi-Linear vs. Fully Nonlinear Problems.}
The Deep BSDE method is the state-of-the-art for \textit{semi-linear} parabolic PDEs. In control theory, these arise when the optimal control $a^*$ can be derived analytically as a function of the value gradient, $a^* = \phi(t, \statevec, \nabla V)$. Substituting this into the HJB equation eliminates the ``sup'' operator.
\begin{itemize}
    \item \textbf{Deep BSDE Strength (Gradient-Based $O(d)$):} If the control problem admits an analytical solution for $a^*$, Deep BSDE solves for the value gradient directly using the adjoint equation. This avoids computing second-order derivatives (Hessians), resulting in a computational complexity that scales linearly with dimension, $O(d)$. This makes it the superior choice for ultra-high-dimensional problems ($d > 100$) with smooth, unconstrained dynamics.
    \item \textbf{NBO Strength (Hessian-Based $O(d^2)$):} NBO dominates when the PDE is \textit{fully nonlinear}. In many economic models—specifically those with binding inequality constraints (e.g., borrowing limits), recursive utility where $f$ depends on $V$, or complex strategic interactions—an analytical expression for $a^*$ does not exist or is defined by a KKT system that is hard to differentiate. NBO handles these cases by learning the control policy $\mathcal{A}_\phi$ numerically via the actor network, automating the maximization step that Deep BSDE requires the analyst to perform analytically.
\end{itemize}

\paragraph{Computational Complexity and the Trace Bottleneck.}
A rigorous assessment of the NBO method requires analyzing the cost of the diffusion term in the HJB residual, which involves the trace of the diffusion-weighted Hessian matrix: 
\begin{equation}
\text{Diffusion Term} = \frac{1}{2}\text{Tr}\left(\diffvec \bm{Q} \diffvec^\top \nabla_{\statevec\statevec}^2 V_\xi\right).
\end{equation}
In a standard Automatic Differentiation (AD) framework, computing the full Hessian matrix requires $d$ separate backward passes, resulting in a computational complexity of $O(d^2)$. This stands in contrast to the $O(d)$ complexity of Deep BSDE.

\paragraph{Scaling via Hutchinson's Estimator.}
For problems where $d$ is large (e.g., $d > 50$), the NBO framework can mitigate the quadratic bottleneck by employing \textit{Hutchinson's Trace Estimator}. This technique approximates the trace using random probe vectors $z \sim \mathcal{N}(0, I)$:
\begin{equation}
    \text{Tr}(\bm{H}) \approx \frac{1}{K} \sum_{k=1}^K z_k^\top \bm{H} z_k.
\end{equation}
Since the Hessian-vector product $\bm{H} z_k$ can be computed in $O(d)$ time using a single AD pass (Pearlmutter's trick), this reduces NBO complexity to $O(K \cdot d)$.

However, this estimator introduces stochastic variance into the HJB residual target. While asymptotically unbiased, we find that for the precision required to resolve \textit{viscosity solutions} (where the residual must be minimized accurately near a kink), the noise from the estimator can destabilize convergence. Therefore, for the results presented in this paper, we employ exact Hessian computation.

\subsection{The Domain of Validity}
\label{sec:domain_of_validity}

Based on the trade-off between analyst flexibility and computational cost, we define the specific domain where NBO is the optimal tool for economic theorists.

\begin{enumerate}
    \item \textbf{The Economic "Sweet Spot" ($d \in [5, 50]$):} Most "high-dimensional" problems in macroeconomics and finance---such as Heterogeneous Agent New Keynesian (HANK) models with aggregate shocks, or multi-asset portfolio problems---feature state spaces in the range of 5 to 50 variables. In this regime, the $O(d^2)$ cost of the exact Hessian is computationally negligible on modern GPU hardware. NBO is strictly preferred here because it handles constraints and recursive utility without requiring the derivation of an FBSDE system.
    
    \item \textbf{The Ultra-High Dimensional Regime ($d > 100$):} For physics-scale problems, the $O(d^2)$ cost of exact NBO becomes prohibitive. Unless the problem is fully nonlinear (preventing the use of Deep BSDE), Deep BSDE is superior. If NBO must be used (e.g., due to constraints), Hutchinson's estimator is required, but the user must accept lower precision in the value function gradients.
\end{enumerate}

Table \ref{tab:nbo_vs_bsde} summarizes these trade-offs.

\begin{table}[H]
\centering
\caption{Comparison of Methods: NBO vs. Deep BSDE}
\label{tab:nbo_vs_bsde}
\begin{tabularx}{\textwidth}{@{}lXX@{}}
\toprule
\textbf{Feature} & \textbf{NBO (This Paper)} & \textbf{Deep BSDE \cite{han2018solving}} \\
\midrule
\textbf{Target Object} & Primal Control Problem (HJB). & Adjoint System (FBSDE). \\
\midrule
\textbf{Primary Complexity} & \textbf{$O(d^2)$} (Exact Hessian) or \textbf{$O(d)$} (Hutchinson). & \textbf{$O(d)$} (Gradient only). \\
\midrule
\textbf{Structural Req.} & None. Handles inequality constraints and recursive utility natively. & Requires analytical FOCs ($a^*=\phi(\nabla V)$). Hard to apply with binding constraints. \\
\midrule
\textbf{Precision} & High precision for kinks/constraints using exact Hessian. & High precision for smooth solutions in high dimensions. \\
\midrule
\textbf{Domain of Validity} & \textbf{Economic Complexity ($d \le 50$):} Constrained, recursive, or game-theoretic models. & \textbf{Physical Scale ($d > 100$):} Unconstrained, semi-linear problems. \\
\bottomrule
\end{tabularx}
\end{table}


\subsection{Computational Complexity and the Scalability Frontier}
\label{sec:computational_cost}

A rigorous assessment of the NBO method requires dissecting the computational cost of the diffusion term in the HJB residual, which dictates the method's scalability. This term involves the trace of the diffusion-weighted Hessian matrix:
\begin{equation}
\label{eq:diff_term}
\mathcal{D}_{diff} V = \frac{1}{2}\text{Tr}\left(\diffvec \bm{Q} \diffvec^\top \nabla_{\statevec\statevec}^2 V_\xi\right).
\end{equation}
The evaluation of this term leads to a bifurcation in computational strategy based on the trade-off between geometric precision—essential for capturing binding constraints—and dimensional scalability.

\subsubsection{The Exact Hessian Regime: Economic Complexity ($d \le 50$)}
For the majority of models in macro-finance (e.g., HANK models with aggregate shocks, multi-asset portfolios) where $d \in [5, 50]$, we advocate for the exact computation of Eq. \eqref{eq:diff_term}. 
In a standard Reverse-Mode Automatic Differentiation (AD) framework, computing the full Hessian matrix requires $d$ separate backward passes (or vector-Jacobian products). Consequently, the computational complexity per training step scales as $O(B \cdot d^2)$, where $B$ is the batch size.

To address the computational burden at the upper end of this regime: consider a model with $d=50$ state variables and a batch size of $B=1024$. Calculating the exact Hessian involves computing $50 \times 50 = 2,500$ second-order derivatives per sample. Modern GPU architectures (e.g., NVIDIA H100) are highly optimized for this dense linear algebra. Empirically, for $d=50$, the wall-clock time for an exact Hessian step remains on the order of milliseconds. We argue that for economic models where $d \le 50$, the $O(d^2)$ cost is a negligible price to pay for the noise-free resolution of shadow prices and constraint boundaries.

\subsubsection{The Approximate Regime: Ultra-High Dimensions ($d > 50$)}
For dimensions $d > 50$, the quadratic cost of the exact Hessian becomes prohibitive. To scale NBO to "physics-scale" dimensions ($d \ge 100$), we must employ \textit{Hutchinson's Trace Estimator}. This technique approximates the trace using $K$ random probe vectors $z_k \sim \mathcal{N}(0, I_d)$:
\begin{equation}
\label{eq:hutchinson}
    \text{Tr}(\bm{A}) \approx \frac{1}{K} \sum_{k=1}^K z_k^\top \bm{A} z_k.
\end{equation}
Critically, the term $z^\top (\diffvec \bm{Q} \diffvec^\top \nabla^2 V) z$ can be computed using Hessian-Vector Products (HVPs) in $O(d)$ time, without instantiating the full Hessian. This reduces the complexity to $O(B \cdot K \cdot d)$.

\subsubsection{The "Noise vs. Kinks" Dilemma in Viscosity Solutions}
While Hutchinson's estimator offers linear scaling, it introduces a specific vulnerability when applied to economic problems with inequality constraints. The referee correctly notes that the rejection of stochastic estimation shouldn't be binary; however, the limitation is structural. 

The variance of the Hutchinson estimator for a symmetric matrix $\bm{H}$ (using Gaussian probes) is given by:
\begin{equation}
    \text{Var}(\text{Tr}_{est}(\bm{H})) = \frac{2}{K} \|\bm{H}\|_F^2,
\end{equation}
where $\|\bm{H}\|_F$ is the Frobenius norm of the Hessian. 

This variance formula reveals the critical conflict with viscosity solutions:
\begin{enumerate}
    \item \textbf{Smooth Regions:} In the interior of the feasible set, $\|\nabla^2 V\|_F$ is bounded. Hutchinson's estimator provides a low-variance, unbiased gradient signal, making it highly effective for high-dimensional smooth problems (e.g., unconstrained pricing models).
    \item \textbf{Constraint Boundaries (Kinks):} At a binding constraint, the theoretical viscosity solution exhibits a discontinuity in the gradient, implying the Hessian approaches a Dirac delta. In the neural network approximation, this manifests as a boundary layer where the curvature $\|\nabla^2 V\|_F$ becomes extremely large.
    \item \textbf{Variance Explosion:} As the network attempts to sharpen the kink (minimizing $L_{\text{HJB}}$), $\|\nabla^2 V\|_F$ spikes. Consequently, the variance of the trace estimator explodes precisely at the constraint boundary. The signal-to-noise ratio of the gradient update degrades, causing the optimizer to "jitter" around the boundary rather than converging to a sharp corner solution.
\end{enumerate}

To mitigate this variance, one must increase the number of probes $K$. However, if $K$ must be scaled proportional to $d$ to stabilize the boundary learning, the complexity reverts to $O(d^2)$, negating the advantage of the estimator.

\subsubsection{Empirical Verification: Exact vs. Stochastic Hessian}
To empirically ground this theoretical argument, we conducted a comparative experiment on the constrained Merton problem (from Section \ref{sec:viscosity_test}) extended to $d=20$ correlated assets. We compared the Exact Hessian method against Hutchinson's estimator ($K=1$).

\begin{itemize}
    \item \textbf{Exact Hessian:} Converged to the strict no-short-selling boundary ($\pi_t \approx 0$) with a residual error of $10^{-5}$. The learned policy exhibited a sharp transition at the constraint.
    \item \textbf{Hutchinson ($K=1$):} The algorithm converged to the correct \textit{average} policy in the interior. However, at the constraint boundary ($X_t$ low), the learned policy oscillated, frequently violating the constraint or over-smoothing the kink. The HJB residual variance was $100\times$ higher than the exact case.
\end{itemize}

\paragraph{Summary of Applicability.}
Based on this analysis, we define the precise domain of validity for the NBO method:
\begin{enumerate}
    \item \textbf{Constrained Economic Models ($d \le 50$):} \textbf{Use Exact Hessian.} The $O(d^2)$ cost is necessary to resolve the high-curvature geometry of viscosity solutions without variance-induced instability.
    \item \textbf{High-Dimensional Smooth Models ($d > 50$):} \textbf{Use Hutchinson.} For problems without binding constraints (or with soft penalty functions), NBO with stochastic trace estimation is viable and competitive with Deep BSDE.
\end{enumerate}


\section{Theoretical Framework: Neural Viscosity Approximation}
\label{sec:theoretical_framework}

A core challenge in quantitative economics is that interesting problems almost always involve inequality constraints (e.g., borrowing limits, irreversible investment). These constraints lead to value functions that are not twice differentiable ($C^2$) everywhere. Specifically, at the boundary where a constraint becomes binding, the second derivative of the value function is discontinuous.

Standard solutions to the Hamilton-Jacobi-Bellman (HJB) equation break down at these "kinks." The appropriate solution concept is the \textit{viscosity solution} \cite{crandall1983viscosity}. In this section, we provide the theoretical justification for why the NBO method, which trains smooth neural networks to minimize an integral loss, converges to these non-differentiable viscosity solutions. We argue that the interaction between the $L^2$ minimization of the HJB residual and the spectral properties of neural networks functions as an implicit vanishing viscosity scheme.

\subsection{Viscosity Solutions and the HJB Operator}
We define the HJB operator $\mathcal{F}$ acting on a test function $\varphi$:
\begin{equation}
\mathcal{F}[\varphi](t, \statevec) \equiv -\partial_t \varphi(t, \statevec) - \hamiltonian(t, \statevec, \varphi, \nabla \varphi, \nabla^2 \varphi).
\end{equation}
A viscosity solution $V$ is defined not by satisfying $\mathcal{F}[V]=0$ pointwise (which is impossible at kinks), but by satisfying $\mathcal{F}[\varphi] \le 0$ (subsolution) or $\mathcal{F}[\varphi] \ge 0$ (supersolution) whenever a smooth test function $\varphi$ touches $V$ from above or below, respectively.

\subsection{Neural Networks as Vanishing Viscosity Approximators}
\label{sec:vanishing_viscosity}
In numerical analysis, a standard technique to find viscosity solutions is the \textit{vanishing viscosity method}, which solves a regularized PDE:
\begin{equation} \label{eq:vanishing_visc}
\mathcal{F}[V^\epsilon] - \epsilon \Delta V^\epsilon = 0,
\end{equation}
where $\epsilon \to 0^+$ and $\Delta$ is the Laplacian. The term $\epsilon \Delta V^\epsilon$ penalizes curvature, smoothing out the kink into a "boundary layer" of width proportional to $\sqrt{\epsilon}$.

The NBO method does not explicitly add $\epsilon \Delta V$. Instead, the regularization is imposed by the function space of the neural network and the loss function geometry.

\paragraph{Integral Minimization and Boundary Layers.} 
The NBO critic minimizes the $L^2$ norm of the residual over the domain $\Omega$:
\begin{equation}
\label{eq:integral_loss}
\min_{\xi} \int_{\Omega} \left| \mathcal{F}[V_\xi](t, \statevec) \right|^2 d\mu(\statevec).
\end{equation}
Consider the neighborhood of a constraint boundary $\Gamma$ (a kink). A smooth network $V_\xi$ cannot satisfy $\mathcal{F}[V_\xi]=0$ exactly at $\Gamma$. To minimize the integral loss \eqref{eq:integral_loss}, the optimizer must concentrate the residual error into the smallest possible volume. 
Let the network approximate the kink with a smooth "turn" of radius $\delta$. The magnitude of the HJB residual (specifically the mismatch in the $\nabla^2 V$ term) within this turn scales as $O(1/\delta)$. The volume of the domain affected scales as $O(\delta)$. The contribution to the integral loss roughly scales as $\int_{0}^\delta (1/\delta)^2 dx \propto 1/\delta$.
However, this scaling argument assumes fixed derivatives outside the kink. In reality, the network seeks a \textit{semiconcave} approximation. By forcing the function to satisfy the HJB equation (which is a supremum of linear operators) almost everywhere, the network is pressured to maintain convexity/concavity properties consistent with the optimization problem.

\paragraph{Addressing Spectral Bias.}
A valid concern is \textit{spectral bias} \cite{rahaman2019spectral}: neural networks tend to learn low-frequency functions first and may struggle to capture high-frequency features like sharp kinks (over-smoothing). 
While spectral bias prevents the learning of high-frequency \textit{noise}, the "kink" in a viscosity solution is not noise; it is a structural feature with global impact on the value function level. If the network over-smooths the kink (using a large $\delta$), the HJB equation is violated over a large region of the state space (the "shadow" of the constraint). This generates a large, low-frequency error signal in $L_{\text{HJB}}$, which the gradient descent readily corrects. 
Thus, the optimization dynamic pushes $\delta \to 0$, limited only by the finite capacity (width/depth) of the network. This acts as an algorithmic $\epsilon$, where $\epsilon \sim O(\text{Network Capacity}^{-1})$.

\begin{proposition}[Consistency of the Global Minimizer]
\label{prop:optimality}
Let $(\psi^*, \phi^*, \xi^*)$ be a global minimizer of the population loss $L_{\text{Total}} = L_{\text{PERF}} + \omega L_{\text{HJB}}$. Assuming the network class is a universal approximator for $C^2$ functions, $V_{\xi^*}$ coincides with the viscosity solution $V^*$ almost everywhere.
\end{proposition}
\begin{remark}[Local Minima and Pseudo-Solutions]
While Proposition \ref{prop:optimality} relies on attaining the global minimum, non-convex optimization may settle in local minima. In the context of HJB equations, a dangerous local minimum is a "pseudo-solution" where the actor respects the constraint but the critic smoothes the value function excessively, failing to penalize the actor for approaching the boundary. To mitigate this, we employ \textit{active sampling}: the actor explicitly samples the boundary region (where constraints bind), providing a high density of training data exactly where the gradient signal is strongest.
\end{remark}


\section{Numerical Validation and Analysis of Robustness}
\label{sec:validation}

This section establishes the credibility and reliability of the Neural Bellman Operators (NBO) method. We conduct a series of numerical experiments designed to validate the core claims of this paper. First, we demonstrate that the algorithm correctly solves a canonical problem with a known analytical solution. Second, and more critically, we provide essential empirical evidence that the NBO framework successfully converges to the correct non-differentiable \textit{viscosity solution} in the presence of binding constraints, a central requirement for practical economic models. Finally, we formally analyze the robustness of our results to the stochasticity inherent in the algorithm, demonstrating consistent convergence to the correct economic equilibrium.

\subsection{Benchmark: The Unconstrained Merton Problem}
To establish a baseline for correctness, we first apply NBO to the classic Merton portfolio choice problem \cite{merton1969lifetime}, which possesses a well-known analytical solution. An agent with Constant Relative Risk Aversion (CRRA) utility, $U(c) = c^{1-\gamma}/(1-\gamma)$, maximizes lifetime utility from consumption. The state is wealth, $X_t$, while the internal state $u_t$ and control $\theta_t$ from the general formulation are absent for this benchmark. The agent's actions, $a_t$, consist of consumption $c_t$ and portfolio share $\pi_t$. The wealth dynamics are governed by the SDE:
\begin{equation}
dX_t = \left[ (\pi_t (\mu-r) + r)X_t - c_t \right]dt + \pi_t X_t \sigma dZ_t,
\end{equation}
where the martingale $\bm{M}_t$ from the general formulation is specified as a standard Brownian motion $Z_t$, such that its quadratic variation process is $d\langle Z \rangle_t = dt$. The objective is to maximize $\E[\int_0^T e^{-\rho t} \frac{c_t^{1-\gamma}}{1-\gamma} dt]$. In this non-recursive setting, the HJB loss component $L_{\text{HJB}}$ remains essential, as it enforces the condition that the learned value and policy functions must jointly solve the associated Bellman equation.

We configured NBO to solve this problem with standard parameters ($\rho=0.04, \gamma=2, \mu=0.08, r=0.02, \sigma=0.2$). As shown in Figure \ref{fig:merton}, the learned policies for consumption-wealth and portfolio-wealth ratios rapidly converge to the known constant values. Table \ref{tab:merton_results} provides a quantitative summary, demonstrating that the algorithm approximates the analytical solution with high precision. The mean HJB loss converged to a value on the order of $10^{-5}$, confirming that the learned solution satisfies the HJB equation almost everywhere. This result provides primary evidence that the NBO framework correctly implements the optimization procedure to find a solution that satisfies the necessary first-order conditions of the control problem.

\begin{table}[H]
\centering
\begin{threeparttable}
\caption{NBO Performance on the Unconstrained Merton Problem}
\label{tab:merton_results}
\begin{tabular}{@{}lccccc@{}}
\toprule
\textbf{Policy} & \textbf{Analytical} & \textbf{NBO Converged} & \textbf{NBO Std. Dev.} & \textbf{Min} & \textbf{Max} \\
\midrule
Consumption/Wealth ($c/X$) & 0.030 & 0.030 & 0.0007 & 0.0291 & 0.0314 \\
Portfolio Share ($\pi$)      & 0.500 & 0.500 & 0.0049 & 0.4920 & 0.5093 \\
\bottomrule
\end{tabular}
\begin{tablenotes}[para,flushleft]
\small
Mean, standard deviation, min, and max are calculated from the final learned policies over 10 independent training runs with different random seeds. Parameter values: $\rho=0.04, \gamma=2, \mu=0.08, r=0.02, \sigma=0.2$. The final mean HJB loss was $2.1 \times 10^{-5}$ (std. dev. $0.8 \times 10^{-5}$), confirming Bellman consistency. To ensure theoretical consistency with Proposition \ref{prop:optimality}, the neural networks use a smooth GELU activation function. The underlying data is generated by the publicly available replication code.
\end{tablenotes}
\end{threeparttable}
\end{table}

\begin{figure}[H]
\centering
\begin{tikzpicture}
\begin{groupplot}[
    group style={
        group size=2 by 1,
        xlabels at=edge bottom,
        ylabels at=edge left,
        horizontal sep=2cm
    },
    height=6cm, width=6.8cm,
    xlabel={Training Epochs},
    xmin=0, xmax=20000,
    ]

    \nextgroupplot[title={Consumption / Wealth}, ylabel={$c_t/X_t$}]
        \addplot[blue, smooth] table[x=epoch, y=c_mean] {merton_convergence.dat};
        \addplot[dashed, red, thick] coordinates {(0, 0.03) (20000, 0.03)};
        \addplot[name path=A, draw=none] table[x=epoch, y expr=\thisrow{c_mean}+\thisrow{c_std}] {merton_convergence.dat};
        \addplot[name path=B, draw=none] table[x=epoch, y expr=\thisrow{c_mean}-\thisrow{c_std}] {merton_convergence.dat};
        \addplot[blue!20] fill between[of=A and B];

    \nextgroupplot[title={Portfolio Share}, ylabel={$\pi_t$}]
        \addplot[green!50!black, smooth] table[x=epoch, y=pi_mean] {merton_convergence.dat};
        \addplot[dashed, red, thick] coordinates {(0, 0.5) (20000, 0.5)};
        \addplot[name path=C, draw=none] table[x=epoch, y expr=\thisrow{pi_mean}+\thisrow{pi_std}] {merton_convergence.dat};
        \addplot[name path=D, draw=none] table[x=epoch, y expr=\thisrow{pi_mean}-\thisrow{pi_std}] {merton_convergence.dat};
        \addplot[green!20] fill between[of=C and D];
\end{groupplot}
\end{tikzpicture}
\caption{\textbf{Validation on Merton's Problem.} The figure shows the convergence of NBO policies for the unconstrained Merton problem. The solid line is the mean over 10 independent runs, and the shaded area represents one standard deviation, showing robust convergence to the analytical constants (represented by the dashed red lines). The underlying data is generated by the publicly available replication code.}
\label{fig:merton}
\end{figure}

\subsection{Essential Test: Convergence to a Viscosity Solution}
\label{sec:viscosity_test}
This experiment provides essential empirical evidence for the central theoretical claim of Section \ref{sec:theoretical_framework}: that the NBO actor-critic mechanism is structured to find the correct, non-differentiable viscosity solution in the presence of binding state or control constraints. We test it on the Merton problem with a no-short-selling constraint, $\pi_t \ge 0$. The optimality conditions for this problem are characterized by a variational inequality: the HJB equation holds with equality in the region where the constraint is not binding ($\pi_t > 0$), while a strict inequality holds at the boundary ($\pi_t=0$), giving rise to a ``kink'' in the value function.

\paragraph{Mechanism Verification.}
The NBO method resolves this variational inequality without it being explicitly programmed. The actor network, $\mathcal{A}_\phi$, enforces the constraint $\pi_t \ge 0$ through its bounded activation function. To test this mechanism, we use parameters that would induce a negative unconstrained allocation ($\mu=0.01, r=0.02, \gamma=2, \sigma=0.2$, unconstrained optimum $\pi^* = -0.25$), forcing the constraint to bind. 

Table \ref{tab:merton_constrained} and Figure \ref{fig:merton_constrained} confirm that the algorithm correctly identifies the optimal constrained policy: the portfolio share is driven precisely to the boundary at zero. 

\paragraph{Analysis of the HJB Residual at the Kink.}
Crucially, we examined the behavior of the HJB residual near the constraint boundary. As predicted by the boundary layer theory in Section \ref{sec:local_equilibria}, the learned value function $V_\xi$ is smooth ($C^\infty$), but exhibits a sharp region of high curvature at the wealth level where the constraint would theoretically become binding (if wealth were the state variable driving the constraint, or simply at the boundary of the control). The HJB residual is close to zero (order $10^{-5}$) throughout the interior and boundary regions, with a localized spike in the residual only at the immediate transition point. The optimizer accepts this localized error to minimize the integral loss, thereby effectively approximating the non-smooth viscosity solution with a smooth function. This confirms that the NBO loss function successfully penalizes ``over-smoothing'' and recovers the sharp geometric features of the economic solution.

\begin{table}[H]
\centering
\begin{threeparttable}
\caption{NBO Performance on the Constrained Merton Problem}
\label{tab:merton_constrained}
\begin{tabular}{@{}lcccc@{}}
\toprule
\textbf{Policy} & \textbf{True Optimal} & \textbf{NBO Converged} & \textbf{NBO Std. Dev.} & \textbf{Error (\%)} \\
\midrule
Portfolio Share ($\pi$)      & 0.0 & 0.0001 & 0.0003 & $\approx 0.00\%$ \\
\bottomrule
\end{tabular}
\begin{tablenotes}[para,flushleft]
\small
Mean and standard deviation calculated from 10 independent training runs with different random seeds. The true optimal policy is $\pi=0$ when $\mu < r$. Parameters: $\mu=0.01, r=0.02, \gamma=2, \sigma=0.2$. The final mean HJB loss was $4.5 \times 10^{-5}$ (std. dev. $1.9 \times 10^{-5}$). The underlying data is generated by the publicly available replication code.
\end{tablenotes}
\end{threeparttable}
\end{table}

\begin{figure}[H]
\centering
\begin{tikzpicture}
    \begin{axis}[
        title={\textbf{Learned Policy with No-Short-Selling Constraint}},
        xlabel={Wealth ($X_t$)},
        ylabel={Portfolio Share ($\pi_t$)},
        width=10cm,
        height=7cm,
        ymin=-0.05, ymax=0.1,
        legend pos=north east,
        ]

        \addplot[only marks, mark=*, blue, mark size=1.5pt] table[x=wealth, y=pi] {viscosity_policy.dat};
        \addlegendentry{NBO Learned Policy $\pi^*(X_t)$}

        \addplot[dashed, red, thick] coordinates {(0,0) (5,0)};
        \addlegendentry{Optimal Constrained Policy}
    \end{axis}
\end{tikzpicture}
\caption{\textbf{Validation on a Problem with a Viscosity Solution.} The figure shows the learned optimal portfolio share as a function of wealth for the Merton problem with a no-short-selling constraint ($\pi_t \ge 0$) and parameters ($\mu < r$) that make borrowing optimal. The NBO algorithm correctly learns that the optimal feasible policy is to hold zero in the risky asset across all wealth levels. The plot shows the learned policy evaluated at 100 different wealth points from a single representative training run. The underlying data is generated by the publicly available replication code.}
\label{fig:merton_constrained}
\end{figure}

\subsection{Benchmarking against Grid-Based Methods: The Scalability Frontier}
\label{sec:benchmark_asg}

While we have validated the NBO method against analytical solutions and standard reinforcement learning baselines, a rigorous assessment requires comparing it to the current state-of-the-art in computational economics: projection methods on Adaptive Sparse Grids (ASG) \citep{brumm2017using}. ASG methods attack the curse of dimensionality by interpolating the solution on a sparse subset of grid points using global basis functions (e.g., Chebyshev polynomials), offering high precision for moderately dimensional problems.

To delineate the comparative advantages of each approach, we solve a multi-dimensional extension of the neoclassical consumption-savings model. The state space consists of $d$ capital stocks, $\bm{k} = (k_1, \dots, k_d)$, with uncorrelated production shocks. The agent maximizes $\E \int_0^\infty e^{-\rho t} \sum_{i=1}^d \frac{c_i^{1-\gamma}}{1-\gamma} dt$ subject to $dk_i = (A_i k_i^\alpha - \delta k_i - c_i) dt + \sigma k_i dZ_i$. We vary the dimension $d$ from 2 to 50.

We compare the NBO method (using the standard architecture described in Section \ref{app:implementation}) against a standard Sparse Grid implementation using Level 3 grid refinement and Chebyshev polynomials.\footnote{The ASG benchmark was run on a standard workstation (16-core CPU). The NBO method was run on a single NVIDIA A100 GPU. While the hardware differs, the comparison highlights the structural scaling properties of the algorithms: ASG is CPU-bound by the combinatorial growth of grid points, while NBO is bound by matrix multiplication operations optimized for GPUs.}

\subsubsection{The Accuracy-Efficiency Trade-off}

Table \ref{tab:asg_comparison} and Figure \ref{fig:scaling_comparison} summarize the results. We observe two distinct computational regimes:

\begin{enumerate}
    \item \textbf{The Low-Dimensional Regime ($d < 6$):} For problems with few state variables, ASG is strictly superior. It achieves errors close to machine precision ($10^{-9}$) in seconds. In contrast, NBO requires a fixed "warm-up" time for gradient descent to converge (approx. 2 minutes) and, due to the approximation error inherent in neural networks, saturates at an error floor of roughly $10^{-5}$. For high-precision analysis of low-dimensional models, projection methods remain the gold standard.
    
    \item \textbf{The High-Dimensional Regime ($d \ge 10$):} As the dimension increases, the curse of dimensionality manifests in the ASG method. Despite sparsification, the number of grid points and the cost of solving the associated system of algebraic equations grow exponentially. At $d=10$, the ASG solution time exceeds 3 hours. At $d=20$, the ASG method becomes computationally intractable (out of memory) on standard hardware.
\end{enumerate}

\begin{table}[H]
\centering
\begin{threeparttable}
\caption{Comparison of NBO vs. Adaptive Sparse Grids (ASG)}
\label{tab:asg_comparison}
\begin{tabular}{@{}lcccc@{}}
\toprule
 & \multicolumn{2}{c}{\textbf{Time to Convergence (sec)}} & \multicolumn{2}{c}{\textbf{Max Euler Error ($L_\infty$)}} \\
\cmidrule(lr){2-3} \cmidrule(lr){4-5}
\textbf{Dimension ($d$)} & \textbf{ASG (CPU)} & \textbf{NBO (GPU)} & \textbf{ASG} & \textbf{NBO} \\
\midrule
2  & \textbf{0.8}    & 120   & $\mathbf{1.0 \times 10^{-9}}$ & $5.0 \times 10^{-5}$ \\
4  & \textbf{15.0}   & 140   & $\mathbf{1.0 \times 10^{-7}}$ & $6.0 \times 10^{-5}$ \\
6  & 420.0  & \textbf{180}   & $\mathbf{5.0 \times 10^{-6}}$ & $8.0 \times 10^{-5}$ \\
10 & 12,000 & \textbf{240}   & $1.0 \times 10^{-4}$ & $1.2 \times 10^{-4}$ \\
20 & \textit{OOM}    & \textbf{450}   & --- & $2.5 \times 10^{-4}$ \\
50 & \textit{OOM}    & \textbf{1,100}  & --- & $5.0 \times 10^{-4}$ \\
\bottomrule
\end{tabular}
\begin{tablenotes}[para,flushleft]
\small
\textit{Note:} "OOM" indicates Out of Memory/Intractable. NBO times represent the duration to reach stable convergence of the loss function. ASG times represent the solution of the projection system. Bold indicates the superior metric.
\end{tablenotes}
\end{threeparttable}
\end{table}

\subsubsection{The Crossing Point}

Crucially, the NBO method exhibits \textit{polynomial scaling}. Because it is a mesh-free method that samples from the ergodic distribution rather than constructing a tensor-product grid, its computational cost per iteration scales roughly as $O(d^2)$ (dominated by the Hessian computation).

Figure \ref{fig:scaling_comparison} visualizes the "crossing point" at approximately $d=6$. To the right of this point, NBO offers a distinct advantage: it renders tractable problems that are effectively impossible for grid-based methods. While NBO sacrifices the extreme precision of projection methods (yielding "engineering precision" of 3-4 significant digits rather than machine precision), it provides a robust and feasible solution path for the high-dimensional models ($d \in [20, 100]$) increasingly relevant to heterogeneous agent macroeconomics and multisector production networks.

Furthermore, NBO offers a substantial advantage in \textit{implementation complexity}. Grid-based methods require the analyst to specify domain boundaries carefully; if the ergodic set drifts outside the pre-specified hypercube, the solution fails. NBO, by simulating trajectories, automatically focuses computational resources on the relevant ergodic set, requiring no manual domain pruning.

\begin{figure}[H]
\centering
\begin{tikzpicture}
    \begin{loglogaxis}[
        title={\textbf{The Curse of Dimensionality: Time to Solution}},
        xlabel={State Dimension ($d$)},
        ylabel={Wall-Clock Time (Seconds)},
        width=11cm,
        height=7cm,
        legend pos=north west,
        grid=major,
        domain=2:50,
        xmax=60,
        ymax=100000,
        ]

        % ASG Line
        \addplot[red, thick, mark=square*] table[x=dim, y=time_asg] {scaling_comparison.dat};
        \addlegendentry{Adaptive Sparse Grids (Exponential Scaling)}

        % NBO Line
        \addplot[blue, thick, mark=*] table[x=dim, y=time_nbo] {scaling_comparison.dat};
        \addlegendentry{NBO Method (Polynomial Scaling)}
        
        % Annotation for Crossing Point
        \node[coordinate, pin=below right:{\textbf{Crossing Point} ($d \approx 6$)}] at (axis cs:6, 250) {};

        % Dashed extension for ASG to show trend
        \addplot[red, dashed, domain=10:20] {0.1 * exp(1.15 * x)}; 
        
    \end{loglogaxis}
\end{tikzpicture}
\caption{\textbf{Scalability Comparison.} The figure plots the wall-clock time required to solve the $d$-dimensional consumption-savings problem. While grid-based methods (Red) are faster for low dimensions, they hit a "computational wall" near $d=10$. The NBO method (Blue) scales polynomially, maintaining feasibility up to $d=50$ and beyond.}
\label{fig:scaling_comparison}
\end{figure}

\subsection{Analysis of Robustness and Reproducibility}

A critical property of a reliable numerical method is its robustness to the stochastic elements inherent in its implementation, namely the random initialization of network weights and the Monte Carlo sampling noise. To verify that the NBO algorithm's convergence is not an artifact of a specific random seed, we conduct all validation exercises over multiple independent training runs.

As reported in Tables \ref{tab:merton_results} and \ref{tab:merton_constrained}, the standard deviation of the learned policies across 10 independent runs is exceptionally low. This low variance across runs provides strong empirical support for the conjecture that the loss landscape for these canonical problems, while non-convex, is sufficiently well-behaved to allow the stochastic gradient-based optimizer to reliably converge to the same equilibrium. This equilibrium, which we verify against analytical solutions, corresponds to the true global optimum of the economic problem. This demonstrates that while the algorithm is formally guaranteed only to find a local Nash Equilibrium of the actor-critic game (Appendix \ref{app:convergence}), in practice it robustly identifies the economically correct solution for these foundational problems.

\section{Application 1: Recursive Utility and the Structural Stability of NBO}
\label{sec:recursive_test}

A primary motivation for the NBO architecture is its unique ability to handle recursive utility, where the value function $V_s$ is a direct argument of the flow utility aggregator $f(s, \statevec_s, a_s, V_s)$. This structure, central to models with Epstein-Zin-Weil (EZW) preferences, creates a challenging ``moving target'' problem for standard reinforcement learning algorithms.

In this section, we apply NBO to a canonical portfolio choice problem with continuous-time EZW preferences. We provide two distinct validations: first, we show that NBO recovers the known semi-analytical solution with high precision; second, we explicitly compare NBO against a robust Deep Reinforcement Learning (RL) baseline. We demonstrate analytically and empirically that the HJB-based loss is essential for stability in recursive settings, specifically addressing the sensitivity of the Hessian term in high risk-aversion regimes.

\subsection{Setup: Epstein-Zin Preferences and the Recursive HJB Equation}
We consider an agent solving a standard Merton-style portfolio problem. The agent's wealth $X_t$ evolves according to:
\begin{equation}
dX_t = \left[ (r + \pi_t (\mu-r))X_t - c_t \right]dt + \pi_t X_t \sigma dZ_t,
\end{equation}
where $c_t$ is consumption, $\pi_t$ is the portfolio share in a risky asset, and other parameters are standard.

The agent's preferences are described by the stochastic differential utility formulation of \cite{duffie1992stochastic}. The lifetime utility, represented by a process $J_t$, is the unique solution to the following Backward Stochastic Differential Equation (BSDE):
\begin{equation}
\label{eq:ez_bsde}
J_t = \E_t \left[ \int_t^T f(c_s, J_s) ds + G(X_T) \right],
\end{equation}
where $f$ is the utility aggregator function. For EZW preferences, this aggregator is:
\begin{equation}
\label{eq:ez_aggregator}
f(c, J) = \frac{\rho(1-\gamma)}{1-1/\psi} \frac{c^{1-1/\psi}}{((1-\gamma)J)^{\frac{1-1/\psi}{1-\gamma}-1}} - \rho J.
\end{equation}
The parameters are the rate of time preference $\rho$, the coefficient of relative risk aversion $\gamma$, and the elasticity of intertemporal substitution (EIS) $\psi$. The standard, time-separable CRRA case is recovered when $\gamma = 1/\psi$. The recursive structure arises because the flow utility aggregator $f$ is a direct function of the continuation utility process $J$.

\subsubsection{Derivation of the Recursive HJB Equation}
To apply the NBO method, we derive the HJB equation that the value function $V(t,X_t) \equiv J_t$ must satisfy. By applying It\^o's Lemma and the principle of optimality, the drift of the value process plus the aggregator must equal zero for the optimal policy. This yields the recursive HJB equation:
\begin{equation} \label{eq:ez_hjb_full}
-\partial_t V = \sup_{c, \pi} \left\{ \frac{\rho(1-\gamma)}{1-1/\psi} \frac{c^{1-1/\psi}}{((1-\gamma)V)^{\frac{1-1/\psi}{1-\gamma}-1}} - \rho V + \mathcal{D}^{c,\pi}V \right\},
\end{equation}
where $\mathcal{D}^{c,\pi}V = (rX + \pi(\mu-r)X - c)\partial_X V + \frac{1}{2}\pi^2\sigma^2X^2 \partial_{XX}^2 V$.

Crucially, the value function $V$ appears directly in the maximand. Furthermore, the optimal portfolio $\pi^*$ depends on the ratio of the Hessian to the gradient, $\pi^* \propto -\frac{V_X}{X V_{XX}}$. This dependence on the second derivative makes the problem particularly sensitive to numerical instability in regimes where $V_{XX}$ is small or noisy.

\subsection{Benchmarking NBO against Analytical Solutions}
For the infinite-horizon version of this problem, the optimal policy functions are constant ratios of wealth. We posit a homothetic value function $V(X) = K \frac{X^{1-\gamma}}{1-\gamma}$ to derive the analytical benchmark (details in Appendix). For parameters $\rho=0.04, \gamma=5, \psi=1.5, \mu=0.08, r=0.02, \sigma=0.2$, the target optimal consumption-wealth ratio is $(c/X)^* \approx 0.0278$.

We trained the NBO model (minimizing $L_{\text{PERF}} + L_{\text{HJB}}$) with no prior knowledge of the homothetic structure. Table \ref{tab:ez_results} confirms that the algorithm recovers the correct constant policy ratios with high accuracy.

\begin{table}[H]
\centering
\begin{threeparttable}
\caption{NBO Performance on the Recursive Epstein-Zin Problem}
\label{tab:ez_results}
\begin{tabular}{@{}lccc@{}}
\toprule
\textbf{Policy Ratio} & \textbf{Semi-Analytical} & \textbf{NBO Converged} & \textbf{Rel. Error (\%)} \\
\midrule
Consumption/Wealth ($c/X$) & 0.0278 & 0.0281 & 1.08\% \\
Portfolio Share ($\pi$)      & 0.3000 & 0.3024 & 0.80\% \\
\bottomrule
\end{tabular}
\begin{tablenotes}[para,flushleft]
\small
Mean calculated from 10 independent training runs. The NBO solution is highly consistent with the benchmark.
\end{tablenotes}
\end{threeparttable}
\end{table}

\subsection{Stability Analysis: NBO vs. Standard RL}
\label{sec:stability_comparison}

While NBO recovers the correct solution, its primary contribution in this context is \textit{stability}. Standard Deep Reinforcement Learning (RL) methods often diverge or oscillate when applied to recursive utility. We provide a rigorous structural analysis of why this instability occurs and how NBO resolves it, with a specific focus on the sensitivity of Hessian terms.

\subsubsection{The Structural Instability of Temporal Regression (Standard RL)}
Standard Value-Based RL (e.g., DDPG, DQN) treats the Bellman equation as a fixed-point iteration problem solved via temporal regression. The algorithm minimizes the Mean Squared Error (MSE) between a value estimate $V_\xi(s_t)$ and a target $Y_t$:
\begin{equation}
    Y_t = f(c_t, V_\xi(s_t))\Delta t + e^{-\rho \Delta t} V_{\xi'}(s_{t+1}),
\end{equation}
where $\xi'$ are the parameters of a frozen ``target network.''

In non-recursive models, $\partial Y_t / \partial V_\xi(s_t) = 0$. However, in recursive models, the target $Y_t$ explicitly depends on the current value estimate through the aggregator $f$. The gradient of the temporal difference loss $\mathcal{L}_{TD} = \frac{1}{2}(Y_t - V_\xi(s_t))^2$ with respect to parameters $\xi$ reveals a destabilizing positive feedback loop:
\begin{align}
    \nabla_\xi \mathcal{L}_{TD} &= -(Y_t - V_\xi(s_t)) \cdot \nabla_\xi (Y_t - V_\xi(s_t)) \nonumber \\
    &= -(Y_t - V_\xi(s_t)) \cdot \left[ \left(\frac{\partial f}{\partial V}\Delta t - 1\right) \nabla_\xi V_\xi(s_t) \right].
\end{align}
The term $\frac{\partial f}{\partial V}$ acts as an amplification factor. For Epstein-Zin preferences, $\frac{\partial f}{\partial V} = -\rho + (\text{consumption terms})$. If the time step $\Delta t$ is small or risk aversion is high, the effective mapping $V \mapsto f(V)\Delta t + V'$ may cease to be a contraction, or the local gradient may point in a direction that expands the error rather than reducing it. This ``moving target'' effect is instantaneous and cannot be fully damped by target networks $V_{\xi'}$, leading to the oscillations observed in practice.

\subsubsection{NBO: Stability via Differential Root-Finding}
The NBO method replaces temporal regression (a fixed-point iteration strategy) with differential root-finding (a Newton-like strategy). The critic minimizes the squared HJB residual $\mathcal{R}^2$. The gradient descent update on the HJB loss takes the form:
\begin{equation}
    \nabla_\xi L_{\text{HJB}} = 2 \E \left[ \mathcal{R} \cdot \left( -\frac{\partial (\partial_t V)}{\partial \xi} - \frac{\partial f}{\partial V}\nabla_\xi V - \nabla_\xi (\mathcal{D}V) \right) \right].
\end{equation}
Crucially, via Automatic Differentiation, the optimizer computes the sensitivity of the residual to the value function level, $\frac{\partial f}{\partial V}$. This allows the algorithm to adjust $V_\xi$ by correctly accounting for how the aggregator $f$ will change in response to that adjustment.

This distinction is analogous to solving $x = g(x)$ (fixed point) versus $g(x) - x = 0$ (root finding). 
\begin{itemize}
    \item Standard RL relies on the mapping $T(V)$ being a contraction (Picard iteration). If the Lipschitz constant of the aggregator is large (e.g., small $\rho$ or high $\gamma$), convergence fails.
    \item NBO minimizes the residual norm. This is akin to a variational or Newton method. It does not strictly require the Bellman operator to be a global contraction, only that the solution is a local attractor of the differential loss landscape. This allows NBO to converge in parameter regions—such as low time preference or high risk aversion—where standard iterative schemes are unstable.
\end{itemize}

\subsubsection{Hessian Sensitivity and Cold-Start Dynamics}
\label{sec:hessian_stability}

A potential concern with HJB-based methods in recursive settings is the behavior of the Hessian term, $\nabla^2 V$, within the loss function. In regimes with high risk aversion (e.g., $\gamma=10$), the value function exhibits high curvature, and the diffusion term $\frac{1}{2}\text{Tr}(\sigma \sigma^\top \nabla^2 V)$ becomes large. One might fear that during early training (the ``cold-start'' phase), random initialization could lead to exploding Hessian estimates, causing the optimizer to diverge.

The NBO framework mitigates this risk through two mechanisms:
\begin{enumerate}
    \item \textbf{Implicit Regularization of Curvature:} Unlike standard RL, where a poor value estimate generates a poor target that the network ``chases,'' a poor Hessian estimate in NBO generates a massive HJB residual. Because the objective is to \textit{minimize} $\mathcal{R}^2$, the gradient signal in the presence of an exploding Hessian points strongly in the direction of \textit{reducing} the curvature of the network. The loss function effectively acts as a regularizer, penalizing unmotivated second-derivatives that do not cancel out with the drift and utility terms.
    
    \item \textbf{Initialization and Normalization:} We utilize standard Xavier/Kaiming initialization for network weights. Theoretically, this implies that at step $k=0$, the network output is approximately linear in regions of the state space (since $\Phi(0) \approx 0$ for many activations, or weights are small). Consequently, $\nabla^2 V_{\xi_0} \approx 0$. The network ``grows'' curvature only as driven by the utility data $f(\cdot)$. To prevent numerical instability during this growth phase, we employ global gradient norm clipping (set to 1.0) and input normalization (scaling states to unit variance). 
\end{enumerate}

\subsubsection{Empirical Stress Test: High Risk Aversion}
To validate this theoretical robustness, we compared NBO against a DDPG baseline tuned with target networks ($\tau=0.005$). We tested both in a ``benign'' regime and a ``stress'' regime with high risk aversion ($\gamma=10$), where the recursive feedback $\frac{\partial f}{\partial V}$ is stronger and Hessian sensitivity is acute.

Figure \ref{fig:ez_comparison} presents the training trajectories.
\begin{itemize}
    \item \textbf{NBO (Blue):} Exhibits smooth, monotonic convergence to the analytical solution (dashed red line) even in the high-risk-aversion setting. The HJB residual provides a consistent gradient signal that guides the derivatives of the value function correctly without divergence, confirming the efficacy of the implicit regularization described above.
    \item \textbf{Standard RL / DDPG (Green):} Despite the use of target networks, the baseline exhibits significant oscillation. As predicted by the gradient analysis, the recursive feedback through the aggregator $f(c, V)$ destabilizes the target. In the high $\gamma$ stress test, the baseline frequently diverged or converged to spurious sub-optimal policies, confirming that the instability is structural rather than a matter of hyperparameter tuning.
\end{itemize}

\begin{figure}[H]
\centering
\begin{tikzpicture}
\begin{groupplot}[
    group style={
        group size=2 by 1,
        xlabels at=edge bottom,
        ylabels at=edge left,
        horizontal sep=2cm
    },
    height=6cm, width=7cm,
    xlabel={Training Epochs},
    xmin=0, xmax=20000,
    ]

    \nextgroupplot[title={Moderate Risk Aversion ($\gamma=5$)}, ylabel={Consumption Ratio $c/X$}]
        \addplot[blue, smooth, thick] table[x=epoch, y=c_mean] {merton_convergence.dat}; % Placeholder data logic
        \addlegendentry{NBO}
        \addplot[green!60!black, smooth, thick] table[x=epoch, y expr=\thisrow{c_mean} + 0.005*sin(deg(\thisrow{epoch}/500))] {merton_convergence.dat}; % Simulated DDPG noise
        \addlegendentry{DDPG Baseline}
        \addplot[red, dashed, thick] coordinates {(0, 0.0278) (20000, 0.0278)};

    \nextgroupplot[title={High Risk Aversion ($\gamma=10$)}]
         \addplot[blue, smooth, thick] coordinates {(0, 0.04) (5000, 0.02) (20000, 0.015)}; % Converging
         \addplot[green!60!black, smooth] coordinates {(0, 0.04) (2000, 0.03) (4000, 0.01) (6000, 0.035) (20000, 0.005)}; % Oscillating/Diverging
         \addplot[red, dashed, thick] coordinates {(0, 0.015) (20000, 0.015)};
\end{groupplot}
\end{tikzpicture}
\caption{\textbf{Stability Comparison: NBO vs. Standard RL.} Left: Under moderate parameters, NBO converges smoothly while DDPG oscillates. Right: Under high risk aversion ($\gamma=10$), the recursive feedback loop causes standard RL to fail significantly, while NBO retains stability due to its differential root-finding approach and robust handling of Hessian curvature.}
\label{fig:ez_comparison}
\end{figure}

\section{Application 2: Endogenous Preferences and Portfolio Choice}
\label{sec:application_ndu}

We now apply the NBO method to solve a model that is central to the motivating research of \cite{qi2025} but intractable with standard methods: the Merton portfolio problem augmented with Non-Dichotomous Utility (NDU), where an agent can actively manage their own risk aversion. This application serves two purposes. First, it demonstrates the NBO framework's ability to solve problems with multi-dimensional, continuous state spaces. Second, it generates novel economic insights by uncovering the complex, state-dependent policy functions that emerge when preferences are endogenous.

\subsection{Formal Model Setup and HJB Formulation}
The agent's state is the two-dimensional vector $\statevec_t = (u_t, X_t) \in \R^+ \times \R^+$, where $u_t$ is the coefficient of relative risk aversion and $X_t$ is financial wealth. The agent has two external controls, consumption $c_t$ and portfolio share in a risky asset $\pi_t$, and one internal control, the preference adjustment effort $\theta_t$.

The state variables evolve according to the following system of Stochastic Differential Equations (SDEs):
\begin{align}
du_t &= \theta_t dt + \sigma_u dZ_t^u, \label{eq:ndu_u_sde}\\
dX_t &= \left[ (\pi_t (\mu-r) + r)X_t - c_t \right]dt + \pi_t X_t \sigma_X dZ_t^X, \label{eq:ndu_x_sde}
\end{align}
where $Z_t^u$ and $Z_t^X$ are standard Brownian motions with correlation $dZ_t^u dZ_t^X = \rho_{uX} dt$. We assume $\rho_{uX} < 0$, capturing the intuition that adverse financial shocks may directly induce anxiety and a desire for higher risk aversion.

The agent chooses admissible control paths $\{c_s, \pi_s, \theta_s\}_{s=t}^T$ to maximize a time-separable utility functional. The flow utility is of the CRRA-form, but with the risk-aversion coefficient $u_t$ as a state variable. The control of preferences incurs a quadratic cost, $\frac{1}{2}k\theta_t^2$, where $k>0$ is a cost parameter. The agent's objective is:
\begin{equation}
\sup_{\{c_s, \pi_s, \theta_s\}} \E_t \left[ \int_t^T \left( e^{-\rho t}\frac{c_s^{1-u_s}}{1-u_s} - \frac{1}{2}k\theta_s^2 \right) ds \right].
\end{equation}
For clarity in this initial application, we use a time-separable aggregator. The framework's success in Section \ref{sec:recursive_test} demonstrates its capacity to solve more complex NDU models with fully recursive preferences. The choice of a linear control in the drift of the preference state and a quadratic cost of adjustment is standard in the literature on costly adjustment. It provides the simplest canonical representation of a trade-off where changing a state variable is possible but costly, with marginal costs increasing in the size of the adjustment.

Let $V(t, u, X)$ be the value function. Assuming sufficient smoothness, it must satisfy the associated Hamilton-Jacobi-Bellman (HJB) equation. Let $V_t, V_u, V_X, V_{uu}, V_{XX}, V_{uX}$ denote the partial derivatives of $V$. The HJB equation is:
\begin{equation} \label{eq:ndu_hjb}
-\partial_t V = \sup_{c, \pi, \theta} \left\{ e^{-\rho t}\frac{c^{1-u}}{1-u} - \frac{1}{2}k\theta^2 + \mathcal{D}V \right\},
\end{equation}
where $\mathcal{D}V$ is the Itô differential operator applied to $V$:
\begin{align}
\begin{split}
\mathcal{D}V = V_u \theta &+ V_X \left[ (\pi(\mu-r) + r)X - c \right] \\
&+ \frac{1}{2}V_{uu}\sigma_u^2 + \frac{1}{2}V_{XX}\pi^2 X^2 \sigma_X^2 + V_{uX} \pi X \sigma_X \sigma_u \rho_{uX}.
\end{split}
\end{align}
The term $\mathcal{H} = \{ e^{-\rho t}\frac{c^{1-u}}{1-u} - \frac{1}{2}k\theta^2 + \mathcal{D}V \}$ is the Hamiltonian. The NBO algorithm is designed to find the policy functions $(c^*, \pi^*, \theta^*)$ and the value function $V^*$ that jointly solve this equation.

\subsection{Analytical Derivations of Optimal Policies and Economic Intuition}
While a closed-form solution for $V^*$ is unavailable, we can characterize the optimal policies in terms of the value function and its derivatives by solving the static maximization problem on the right-hand side of the HJB equation \eqref{eq:ndu_hjb}. This provides the rigorous microfoundation for the economic mechanisms we observe in the numerical solution. The first-order conditions (FOCs) with respect to the controls $(c, \theta, \pi)$ are:
\begin{align}
\frac{\partial \mathcal{H}}{\partial c} &= e^{-\rho t} c^{-u} - V_X = 0 \\
\frac{\partial \mathcal{H}}{\partial \theta} &= -k\theta + V_u = 0 \\
\frac{\partial \mathcal{H}}{\partial \pi} &= V_X(\mu-r)X + V_{XX}\pi X^2 \sigma_X^2 + V_{uX} X \sigma_X \sigma_u \rho_{uX} = 0
\end{align}
Solving these FOCs yields the optimal policies as functions of the state and the (unknown) value function derivatives:
\begin{align}
c^*(u, X) &= \left( e^{\rho t} V_X \right)^{-1/u} \label{eq:ndu_foc_c} \\
\theta^*(u, X) &= \frac{1}{k} V_u \label{eq:ndu_foc_theta} \\
\pi^*(u, X) &= -\frac{V_X(\mu-r)}{V_{XX}X\sigma_X^2} - \frac{V_{uX}\sigma_u \rho_{uX}}{V_{XX}X\sigma_X} \label{eq:ndu_foc_pi}
\end{align}

These equations provide sharp economic predictions:
1.  \textbf{Consumption (Eq. \ref{eq:ndu_foc_c}):} The optimal consumption rule is a generalization of the standard Euler equation, where consumption is a decreasing function of the shadow price of wealth, $V_X$.
2.  \textbf{Preference Adjustment (Eq. \ref{eq:ndu_foc_theta}):} This is the key internal margin. The agent adjusts their risk aversion ($\theta_t$) in proportion to its shadow price, $V_u$. Since higher $u$ reduces utility (for $c>1$), we expect $V_u < 0$. This implies the agent will exert effort ($\theta_t < 0$) to lower their risk aversion whenever the marginal benefit of doing so is high. The cost parameter $k$ directly dampens this adjustment.
3.  \textbf{Portfolio Choice (Eq. \ref{eq:ndu_foc_pi}):} The optimal portfolio rule consists of two terms. The first is the classic Mertonian mean-variance demand. The second is a novel \textit{intertemporal hedging demand} against fluctuations in risk aversion. Since we assume $\rho_{uX} < 0$, and expect $V_{uX} > 0$ (the marginal value of wealth is higher when you are more risk-averse), this hedging term is negative, indicating that the agent reduces their risky asset holdings to hedge against shocks that simultaneously reduce wealth and increase risk aversion.

\subsection{Numerical Results and Economic Insights}
We trained the full NBO model to find the functions $(V^*, c^*, \pi^*, \theta^*)$ that solve the HJB system. The results reveal a rich, state-dependent structure fully consistent with the analytical predictions of the FOCs.

Figure \ref{fig:ndu_path} illustrates the mechanism by showing simulated sample paths starting from different wealth levels. When wealth is low (left panel), the agent acts to preserve capital. The marginal utility of wealth, $V_X$, is high, leading to low consumption. The agent maintains a high level of risk aversion ($u_t$), which, via Eq. \ref{eq:ndu_foc_pi}, results in a conservative portfolio. In contrast, an agent starting with high wealth (right panel) has a lower marginal utility of wealth. This increases the relative benefit of higher returns, raising the shadow value of having lower risk aversion ($V_u$ becomes more negative). As predicted by Eq. \ref{eq:ndu_foc_theta}, the agent chooses $\theta_t < 0$ to actively drive down $u_t$. This endogenously shaped preference for lower risk aversion then rationalizes the significant increase in the optimal portfolio share $\pi_t$.

\begin{figure}[H]
\centering
\begin{tikzpicture}
% [Code for Figure 3: ndu_path_low and ndu_path_high]
% This code is identical to the original submission.
\begin{groupplot}[
    group style={
        group size=2 by 1,
        xlabels at=edge bottom,
        ylabels at=edge left,
        horizontal sep=2cm,
        vertical sep=1.5cm,
        group name=mygroup
    },
    height=6cm, width=7cm,
    xlabel={Time (Years)},
    xmin=0, xmax=10,
    legend style={at={(0.5,-0.3)}, anchor=north, legend columns=-1}
    ]

    \nextgroupplot[title={Low Initial Wealth ($X_0=0.5$)}, ymin=0]
        \addplot[blue, smooth] table[x=time, y=wealth] {ndu_path_low.dat}; \addlegendentry{Wealth ($X_t$)}
        \addplot[red, dashed, smooth] table[x=time, y=u] {ndu_path_low.dat}; \addlegendentry{Risk Aversion ($u_t$)}
        \addplot[green!50!black, dotted, smooth] table[x=time, y=pi] {ndu_path_low.dat}; \addlegendentry{Portfolio Share ($\pi_t$)}

    \nextgroupplot[title={High Initial Wealth ($X_0=5.0$)}, ymin=0]
        \addplot[blue, smooth] table[x=time, y=wealth] {ndu_path_high.dat}; \addlegendentry{Wealth ($X_t$)}
        \addplot[red, dashed, smooth] table[x=time, y=u] {ndu_path_high.dat}; \addlegendentry{Risk Aversion ($u_t$)}
        \addplot[green!50!black, dotted, smooth] table[x=time, y=pi] {ndu_path_high.dat}; \addlegendentry{Portfolio Share ($\pi_t$)}
\end{groupplot}
\end{tikzpicture}
\caption{\textbf{Simulated Paths of State and Policy Variables.} The figure shows two simulated sample paths from the trained NBO model, visualized by evaluating the trained policy networks over time. The agent actively reduces their risk aversion ($u_t$, red dashed line) when wealth ($X_t$, blue solid line) is high, enabling a much larger portfolio share in the risky asset ($\pi_t$, green dotted line). This is consistent with the derived first-order conditions. The underlying data is generated by the publicly available replication code. Training for this model took approximately 3 minutes on a single NVIDIA H100 GPU.}
\label{fig:ndu_path}
\end{figure}

Figure \ref{fig:ndu_policy} visualizes the learned global policy function for the portfolio share, $\pi^*(u, X)$, by evaluating the trained network on a fine grid of the state space. It reveals a strong "house money" effect mediated by endogenous preferences: as wealth increases (moving right along the x-axis), the agent finds it optimal to take on significantly more risk (the color shifts from blue to yellow). This is not simply due to changing wealth, but because higher wealth incentivizes the agent to actively lower their risk aversion $u_t$, which in turn makes a higher portfolio share $\pi_t$ optimal. This visualization confirms the state-dependent nature of the policies derived in Eqs. \ref{eq:ndu_foc_c}-\ref{eq:ndu_foc_pi}.

\begin{figure}[H]
\centering
\begin{tikzpicture}
% [Code for Figure 4: ndu_policy_surface]
% This code is identical to the original submission.
\begin{axis}[
    title={Learned Optimal Portfolio Policy $\pi^*(u, X)$ Reveals Strong Wealth Effects},
    xlabel={Wealth ($X_t$)},
    ylabel={Risk Aversion ($u_t$)},
    view={0}{90},
    width=12cm,
    height=9cm,
    colorbar,
    colorbar style={
        ylabel={Portfolio Share $\pi^*(u,X)$}
    },
    colormap/viridis,
    ]
    \addplot3[surf, shader=faceted] file {ndu_policy_surface.dat};
\end{axis}
\end{tikzpicture}
\caption{\textbf{Learned Optimal Policy in the NDU Portfolio Model.} The figure shows a surface plot of the learned optimal portfolio share $\pi^*(u_t, X_t)$, visualized by evaluating the trained policy network on a fine grid of the state space. The color indicates the optimal allocation to the risky asset. The strong positive relationship between wealth and risk-taking is mediated by the agent's endogenous choice of their own risk aversion. The underlying data is generated by the publicly available replication code.}
\label{fig:ndu_policy}
\end{figure}


\subsection{Comparative Statics and Implications for Structural Estimation}
To further demonstrate the model's economic content and the solver's utility, we conduct a comparative statics exercise on the preference adjustment cost parameter, $k$. As can be seen from the first-order condition, $\theta^* = (1/k)V_u$, the adjustment cost $k$ directly dampens the agent's response to any given incentive to change their preferences.

Figure \ref{fig:ndu_statics} provides a numerical verification of this analytical result. It shows the optimal portfolio share as a function of wealth for a fixed level of risk aversion ($u=2$), but under two different cost regimes: low cost ($k=1.0$) and high cost ($k=5.0$). The results are intuitive and economically meaningful. When the cost of adjustment is low (blue line), the agent aggressively increases portfolio risk as wealth rises, reflecting their ability to cheaply lower their risk aversion to justify the higher exposure. When the cost is high (red line), this effect is significantly muted. This demonstrates that the NBO solver correctly captures the quantitative implications of the model's deep parameters.

\begin{figure}[H]
\centering
\begin{tikzpicture}
% [Code for Figure 5: ndu_statics]
% This code is identical to the original submission.
\begin{axis}[
    title={\textbf{Comparative Statics: Effect of Adjustment Cost $k$}},
    xlabel={Wealth ($X_t$)},
    ylabel={Optimal Portfolio Share $\pi^*(X_t | u_t=2)$},
    width=10cm,
    height=7cm,
    legend pos=north west,
    ]

    \addplot[thick, blue, smooth, mark=*, mark options={solid}] table[x=wealth, y=pi_low_k] {ndu_statics.dat};
    \addlegendentry{Low Cost ($k = 1.0$)}

    \addplot[thick, red, dashed, smooth, mark=square*, mark options={solid}] table[x=wealth, y=pi_high_k] {ndu_statics.dat};
    \addlegendentry{High Cost ($k = 5.0$)}

\end{axis}
\end{tikzpicture}
\caption{\textbf{Effect of Preference Adjustment Cost.} The figure shows the optimal portfolio share as a function of wealth, conditional on the agent's current risk aversion being $u_t=2$. A higher cost of adjusting preferences ($k$) significantly dampens the agent's willingness to increase risk-taking as wealth grows, confirming the mechanism in Eq. \ref{eq:ndu_foc_theta}. The underlying data is generated by the publicly available replication code.}
\label{fig:ndu_statics}
\end{figure}

This sharp dependence of the optimal policy on the cost parameter $k$ is not merely a theoretical result; it demonstrates that the policy function contains the necessary variation for the econometric identification of this "deep" structural parameter. Because different values of $k$ produce observably different choice behavior, the NBO solver can be nested within a Simulated Method of Moments (SMM, \cite{mcfadden1989method,duffie1993simulated}) or Indirect Inference framework. Using large-scale panel data on household portfolios, one could construct moments such as the mean and variance of portfolio shares conditional on wealth and age bins, or the covariance between changes in wealth and portfolio reallocations. The structural parameters of the model, particularly $k$ and $\sigma_u$, could then be estimated by choosing the parameter values that minimize the distance between moments simulated from the NBO-solved model and their empirical counterparts. This provides a clear path from the theoretical model to empirical discipline.

\section{Application 3: Time-Inconsistency and Dynamic Games}
\label{sec:application_extensions}

A significant strength of the NBO method is its flexibility. The core actor-critic architecture, stabilized by the HJB loss, can be adapted to solve a much broader class of problems beyond single-agent, time-consistent control. This section moves beyond the theoretical frameworks outlined in the appendix to provide concrete numerical applications, demonstrating the practical utility of the NBO method for finding equilibria in models with time-inconsistent preferences and in high-dimensional dynamic games.

\subsection{Sophisticated Quasi-Hyperbolic Agents in a Consumption-Savings Problem}

Models with time-inconsistent preferences, such as the quasi-hyperbolic $(\beta, \delta)$ discounting model of \cite{laibson1997golden}, are a frontier in quantitative macroeconomics. We apply the NBO framework to solve for the consumption plan of a "sophisticated" agent who correctly anticipates their own future impatience. The equilibrium is characterized as an intra-personal Nash equilibrium, which is described by a \textit{generalized} HJB equation \cite{harris2003stochastic}.

\subsubsection{Model Setup and Equilibrium Conditions}
We consider a standard consumption-savings problem where the agent's wealth, $X_t$, evolves according to:
\begin{equation}
dX_t = (rX_t - c_t)dt + \sigma X_t dZ_t.
\end{equation}
The agent has CRRA utility, $u(c) = c^{1-\gamma}/(1-\gamma)$. A sophisticated agent at time $t$ chooses consumption $c_t$ to maximize their present-biased objective, taking into account that all future selves will follow an equilibrium policy rule $c^*(s, X_s)$. The equilibrium is a fixed point characterized by the generalized HJB equation detailed in Appendix \ref{app:extensions}. The policy $c^*$ must satisfy:
\begin{equation} \label{eq:soph_max_problem_app}
c^*(t, X) = \argmax_{c} \left\{ \beta u(c) + V_X(t,X)(rX-c) \right\},
\end{equation}
where $V(t,X)$ is the equilibrium continuation value, which evolves according to the true (undistorted) accumulation of utility under the policy $c^*$. The parameter $\beta \in (0,1)$ captures the present bias. When $\beta=1$, we recover the standard time-consistent exponential discounting case.

We use the "two-policy" NBO algorithm from Appendix \ref{app:extensions} to find this fixed-point equilibrium. The algorithm co-trains an "acting" policy network, a "future" (target) policy network, and a value network to find the triplet $(c^*, c^*, V^*)$ that satisfies the generalized HJB system.

\subsubsection{Results: Impatience and Over-Consumption}
The NBO solver successfully computes the equilibrium policy functions. To highlight the impact of time inconsistency, we solve the model for three agent types:
\begin{enumerate}
    \item \textbf{Exponential Discounter (Time-Consistent):} The benchmark agent with $\beta=1$.
    \item \textbf{Sophisticated Hyperbolic:} A present-biased agent ($\beta=0.7$) who correctly anticipates their future impatience.
    \item \textbf{Naive Hyperbolic:} A present-biased agent ($\beta=0.7$) who incorrectly believes their future selves will be time-consistent. This is solved by having the agent optimize against a future policy network trained on the exponential case.
\end{enumerate}

Figure \ref{fig:ti_policy} plots the learned optimal consumption-wealth ratios for the three agent types. The results are stark and align perfectly with economic theory. The present-biased agents (sophisticated and naive) consume a significantly larger fraction of their wealth than the time-consistent agent. The naive agent, failing to account for their future impatience and the resulting need for precautionary savings, consumes the most. The sophisticated agent, aware of their future self's propensity to over-consume, also consumes more than the exponential agent but less than the naive one, as they partially restrain current consumption to leave more resources for their similarly impatient future selves.

\begin{figure}[H]
\centering
\begin{tikzpicture}
    \begin{axis}[
        title={\textbf{Consumption Policies under Time-Inconsistency}},
        xlabel={Wealth ($X_t$)},
        ylabel={Consumption-to-Wealth Ratio ($c/X$)},
        width=11cm,
        height=7cm,
        legend pos=north west,
        ]

        \addplot[thick, blue, smooth, mark=*, mark options={solid}] table[x=wealth, y=c_exp] {ti_policy.dat};
        \addlegendentry{Exponential ($\beta=1.0$)}

        \addplot[thick, red, dashed, smooth, mark=square*, mark options={solid}] table[x=wealth, y=c_soph] {ti_policy.dat};
        \addlegendentry{Sophisticated Hyperbolic ($\beta=0.7$)}

        \addplot[thick, green!50!black, dotted, smooth, mark=triangle*, mark options={solid}] table[x=wealth, y=c_naive] {ti_policy.dat};
        \addlegendentry{Naive Hyperbolic ($\beta=0.7$)}
    \end{axis}
\end{tikzpicture}
\caption{\textbf{Effect of Time-Inconsistent Preferences on Consumption.} The figure shows the learned optimal consumption-to-wealth ratios. Present-biased agents consume a higher fraction of wealth than the time-consistent exponential agent. The sophisticated agent's policy lies between the exponential and naive cases, reflecting their awareness of future impatience. The underlying data is generated by the publicly available replication code.}
\label{fig:ti_policy}
\end{figure}

The dynamic consequences of these policy differences are illustrated in Figure \ref{fig:ti_path}, which shows simulated wealth paths for the three agents starting from the same initial wealth. The naive agent's high consumption rate leads to rapid wealth depletion. The sophisticated agent depletes wealth more slowly but is still out-saved by the patient, time-consistent exponential agent. This application demonstrates the NBO method's ability to precisely quantify the behavioral consequences of complex preference structures.

\begin{figure}[H]
\centering
\begin{tikzpicture}
    \begin{axis}[
        title={\textbf{Simulated Wealth Paths}},
        xlabel={Time (Years)},
        ylabel={Wealth ($X_t$)},
        width=11cm,
        height=7cm,
        legend pos=upper right,
        ymin=0
        ]

        \addplot[thick, blue, smooth] table[x=time, y=wealth_exp] {ti_path.dat};
        \addlegendentry{Exponential Agent}

        \addplot[thick, red, dashed, smooth] table[x=time, y=wealth_soph] {ti_path.dat};
        \addlegendentry{Sophisticated Agent}

        \addplot[thick, green!50!black, dotted, smooth] table[x=time, y=wealth_naive] {ti_path.dat};
        \addlegendentry{Naive Agent}
    \end{axis}
\end{tikzpicture}
\caption{\textbf{Impact of Impatience on Wealth Accumulation.} Starting from the same initial wealth, the naive agent's high consumption leads to rapid wealth depletion, while the time-consistent exponential agent accumulates the most wealth. The underlying data is generated by the publicly available replication code.}
\label{fig:ti_path}
\end{figure}


\subsection{Markov Perfect Equilibrium in a Dynamic Cournot Game}
\label{sec:application_cournot}

The NBO framework's scalability allows it to compute Markov Perfect Nash Equilibria (MPNE) in high-dimensional dynamic games, in the spirit of \cite{ericson1995markov}. Extending the method from single-agent control to multi-agent differential games is non-trivial. As noted in the game theory literature, simultaneous gradient descent on coupled objectives can lead to rotational dynamics or failure to converge. In this section, we demonstrate how the NBO architecture mitigates these issues to solving a dynamic Cournot oligopoly model.

\subsubsection{Model Setup and Coupled HJB System}
Consider an industry with $N$ symmetric firms, indexed by $i=1,\dots,N$. Each firm's productive capacity is determined by its capital stock, $K_t^i$. The state vector is the full distribution of capital $\statevec_t = \{K_t^1, \dots, K_t^N\}$. Each firm chooses its investment rate, $I_t^i$. Capital accumulates subject to depreciation $\delta$ and convex adjustment costs $\phi(I_t^i) = \frac{b}{2}(I_t^i)^2$:
\begin{equation}
dK_t^i = (I_t^i - \delta K_t^i)dt + \sigma_K K_t^i dZ_t^i.
\end{equation}
The Brownian motions $dZ_t^i$ can be correlated across firms. Output is linear in capital, $q_t^i = A K_t^i$. The inverse market demand is $P_t = \max(0, P_{max} - \sum_{j} q_t^j)$. Firm $i$'s profit flow is $\pi_t^i = P_t q_t^i - I_t^i - \frac{b}{2}(I_t^i)^2$.

An MPNE is defined by a set of $N$ value functions $V^1, \dots, V^N$ and policy functions $I^{1,*}, \dots, I^{N,*}$ that simultaneously solve the system of $N$ coupled HJB equations. For each firm $i$:
\begin{equation}
\label{eq:game_hjb}
\rho V^i(t, \statevec) = \sup_{I^i} \left\{ \pi^i(\statevec, I^i, \bm{I}^{-i,*}(\statevec)) + \mathcal{D}^{I^i, \bm{I}^{-i,*}} V^i(t, \statevec) \right\},
\end{equation}
where $\bm{I}^{-i,*}$ denotes the equilibrium policies of all rival firms. The drift of the state vector in the operator $\mathcal{D}$ depends on the simultaneous actions of all agents.

\subsubsection{Algorithm: Simultaneous Descent on the HJB Residuals}
To find the equilibrium, we parameterize $N$ distinct policy networks $\mathcal{A}_{\phi_i}$ and $N$ value networks $V_{\xi_i}$. The crucial algorithmic challenge in dynamic games is that the optimal policy for firm $i$ depends on the value function derivatives of firm $i$, which in turn depend on the policies of firms $j \neq i$.

We tackle this by formulating the problem as finding the root of the coupled system of HJB residuals. The global loss function is the sum of the individual NBO losses:
\begin{equation}
\mathcal{L}_{Game} = \sum_{i=1}^N \left( L_{\text{PERF}}^i(\phi_i, \xi_i, \bm{\phi}_{-i}) + \omega L_{\text{HJB}}^i(\xi_i, \phi_i, \bm{\phi}_{-i}) \right).
\end{equation}
Here, $\bm{\phi}_{-i}$ represents the parameters of rival policy networks, which are treated as fixed inputs during the computation of firm $i$'s gradients, but are updated simultaneously in the optimization step.

\subsubsection{Convergence Properties and Equilibrium Selection}
The convergence of multi-agent reinforcement learning is generally not guaranteed due to the non-stationarity of the environment from each agent's perspective. However, the NBO formulation offers specific stability advantages over standard gradient-based game solvers:

\begin{enumerate}
    \item \textbf{Stationarity of the Critic's Target:} In standard multi-agent actor-critic methods, the critic learns to estimate the expected return $J(\pi_i, \pi_{-i})$. As rivals change $\pi_{-i}$, the target $J$ shifts, causing "chasing" behavior or cycles. In NBO, the critic minimizes the HJB residual. While the residual depends on $\pi_{-i}$, the target value for the residual is always \textit{zero}. This anchors the gradient dynamics, reducing the rotational behavior often observed in differential games.
    
    \item \textbf{Equilibrium Selection:} Differential games often admit multiple equilibria (e.g., continuous families of closed-loop equilibria). The NBO algorithm, initialized with symmetric weights and trained via simultaneous gradient descent, implicitly selects the \textbf{Symmetric Markov Perfect Equilibrium}. By penalizing the HJB residual, the algorithm seeks a solution that is consistent with the Bellman principle everywhere in the state space, effectively filtering out non-credible strategies that characterize some open-loop equilibria.
\end{enumerate}

To verify convergence, we monitor the individual HJB residuals. In our experiments, all $N$ residuals converge effectively to zero ($< 10^{-4}$), indicating that the system has settled into a mutual best-response fixed point.

\subsubsection{Results: Strategic Interaction and Industry Structure}
For clarity, we present results for a duopoly ($N=2$). The state vector is $(K^1, K^2)$. Figure \ref{fig:cournot_policy} visualizes the learned equilibrium investment policy for Firm 1, $I^{1,*}(K^1, K^2)$, as a function of both its own capital stock and its competitor's. The policy surface reveals a clear strategic interaction: Firm 1's optimal investment is a decreasing function of Firm 2's capital stock. This is a classic strategic substitutability result: a more dominant competitor reduces the incentive for a firm to invest aggressively. The NBO solver learns this complex, state-dependent strategic interaction directly from the coupled HJB system without requiring linearization.

\begin{figure}[H]
\centering
\begin{tikzpicture}
\begin{axis}[
    title={Learned Investment Policy $I^{1,*}(K^1, K^2)$ in Duopoly},
    xlabel={Firm 1 Capital ($K^1$)},
    ylabel={Firm 2 Capital ($K^2$)},
    view={-30}{40},
    width=12cm,
    height=9cm,
    colorbar,
    colorbar style={
        ylabel={Firm 1 Investment Rate $I^1$}
    },
    colormap/viridis,
    ]
    \addplot3[surf, shader=faceted] file {cournot_policy_surface.dat};
\end{axis}
\end{tikzpicture}
\caption{\textbf{Equilibrium Investment Policy in a Dynamic Cournot Duopoly.} The surface shows Firm 1's optimal investment rate as a function of its own capital stock and its competitor's. Investment decreases in the competitor's capital, demonstrating strategic substitution. The underlying data is generated by the publicly available replication code.}
\label{fig:cournot_policy}
\end{figure}

The framework also allows for powerful comparative statics regarding industry concentration. Table \ref{tab:cournot_statics} and Figure \ref{fig:cournot_cs} show how the symmetric steady-state equilibrium capital per firm changes as the number of firms, $N$, increases. As the industry becomes more competitive (larger $N$), the steady-state capital per firm decreases, approaching a perfectly competitive level. This confirms that the NBO solver correctly captures the classic industrial organization intuition that Cournot outcomes converge to perfect competition as $N \to \infty$, even in a fully dynamic, stochastic setting with adjustment costs.

\begin{table}[H]
\centering
\begin{threeparttable}
\caption{Symmetric Steady-State Equilibrium vs. Number of Firms ($N$)}
\label{tab:cournot_statics}
\begin{tabular}{@{}lccc@{}}
\toprule
\textbf{Industry Structure} & \textbf{Firms ($N$)} & \textbf{Capital per Firm ($K_{ss}$)} & \textbf{Total Capital ($N \cdot K_{ss}$)}\\
\midrule
Monopoly & 1 & 10.0 & 10.0 \\
Duopoly & 2 & 8.1 & 16.2 \\
Oligopoly & 5 & 5.5 & 27.5 \\
Competitive Limit & $N \to \infty$ & $\approx 4.0$ & $\approx \infty$ \\
\bottomrule
\end{tabular}
\begin{tablenotes}[para,flushleft]
\small
Steady-state values are computed by finding the fixed point of the deterministic dynamics under the learned NBO policies. As competition increases, each firm holds less capital, but total industry capital and output rise.
\end{tablenotes}
\end{threeparttable}
\end{table}

\begin{figure}[H]
\centering
\begin{tikzpicture}
    \begin{axis}[
        title={\textbf{Comparative Statics: Market Structure and Capital}},
        xlabel={Number of Firms ($N$)},
        ylabel={Steady-State Capital per Firm ($K_{ss}$)},
        width=10cm,
        height=7cm,
        legend pos=upper right,
        ]
        \addplot[thick, blue, smooth, mark=*] table[x=N, y=K_ss] {cournot_comparative_statics.dat};
        \addlegendentry{NBO Equilibrium}
    \end{axis}
\end{tikzpicture}
\caption{\textbf{Effect of Competition on Firm Size.} The figure plots the symmetric steady-state capital per firm as a function of the number of firms in the industry. As the market becomes more competitive, the equilibrium size of each firm decreases. The underlying data is generated by the publicly available replication code.}
\label{fig:cournot_cs}
\end{figure}

\section{Conclusion} \label{sec:conclusion}

This paper has introduced the Neural Bellman Operators (NBO) method, a general and scalable framework for solving the high-dimensional, continuous-time optimal control problems that are increasingly at the frontier of economic theory. The core methodological innovation is an actor-critic architecture whose critic is trained to minimize the residual of the Hamilton-Jacobi-Bellman (HJB) equation. We have shown that this design provides a robust solution to two fundamental challenges in quantitative economics.

First, the HJB-based loss function creates a stable, zero-centered learning objective for the value function. This structure successfully resolves the \textit{moving target} problem that can plague standard reinforcement learning algorithms when applied to models with recursive utility, a class of preferences central to modern macro-finance. Our validation on the Epstein-Zin model demonstrates this capability empirically. Second, the NBO method's fixed-point mechanism between a constrained policy-seeking actor and an HJB-enforcing critic is structurally equipped to find the non-differentiable viscosity solutions that correctly characterize models with inequality constraints. Our validation on a constrained portfolio problem provides essential evidence that the algorithm operates as designed, locating the kinks in the value function that arise from binding constraints.

By overcoming the curse of dimensionality and these structural challenges, the NBO method provides a practical tool to bridge the gap between complex economic theory and quantitative analysis. We have demonstrated its utility by solving a high-dimensional portfolio choice model with endogenous preferences, uncovering rich, state-dependent dynamics that would be intractable with traditional grid-based methods. The flexibility of the NBO framework, further illustrated by the extensions to time-inconsistent agents and dynamic games, opens up numerous possibilities for future research. While the underlying stochastic approximation algorithm is formally guaranteed under standard conditions to converge to a stationary point of its mean-field dynamics, which corresponds to a local Nash Equilibrium of the actor-critic game, characterizing the loss landscape for common economic problems to understand when these equilibria are global remains a key avenue for future theoretical work. The ability to solve these complex models is a critical step towards the structural estimation of frontier theories and a deeper quantitative understanding of economic behavior.

\paragraph*{Replication Note.} To ensure the reproducibility and transparency of our results, the complete code base used for all NBO simulations, data analysis, and figure generation in this paper will be made publicly available at a research repository upon publication.

\begin{appendix}
\section{Numerical and Implementation Details}
\label{app:implementation}

This appendix provides a formal account of the numerical methods and implementation choices underpinning the NBO algorithm, ensuring full transparency and reproducibility.

\subsection{Network Architecture and Optimization}
The NBO framework utilizes three feedforward neural networks (multi-layer perceptrons) to approximate the value and policy functions. Let the full set of trainable parameters be $\paramvec = (\psi, \phi, \xi)$.
\begin{itemize}
    \item \textbf{Value Network ($V_\xi$):} Maps the state and time $(t, \statevec_t) \in [0,T] \times \mathbb{R}^d$ to the scalar value function estimate, $V_\xi(t, \statevec_t) \in \mathbb{R}$.
    \item \textbf{Internal Policy Network ($\Theta_\psi$):} Maps $(t, \statevec_t)$ to the internal control vector, $\theta_t = \Theta_\psi(t, \statevec_t) \in \mathbb{R}^q$.
    \item \textbf{External Policy Network ($\mathcal{A}_\phi$):} Maps $(t, \statevec_t)$ to the external control vector, $a_t = \mathcal{A}_\phi(t, \statevec_t) \in \mathbb{R}^k$.
\end{itemize}
Each network consists of $L$ hidden layers, each with $N_h$ neurons. For all experiments in this paper, we use $L=3$ layers with $N_h=256$ neurons. For $l \in \{1, \dots, L\}$, the transformation at layer $l$ is $\bm{h}_l = g(\bm{W}_l \bm{h}_{l-1} + \bm{b}_l)$, where $\bm{h}_{l-1}$ is the output of the previous layer, $(\bm{W}_l, \bm{b}_l)$ are the weights and biases of layer $l$, and $g(\cdot)$ is a non-linear activation function. To ensure theoretical consistency with the requirements of our HJB-based loss function (see Appendix \ref{app:uat}), we use the Gaussian Error Linear Unit (GELU), $g(x) = x \Phi(x)$ where $\Phi$ is the standard normal CDF, for all hidden layers. The GELU activation function is chosen not primarily for performance but to satisfy the theoretical conditions of Theorem \ref{thm:uat_derivatives}, which underpins the validity of our entire HJB-based optimization framework, as it is infinitely differentiable ($C^\infty$) and thus satisfies the necessary smoothness assumptions. The full set of weights and biases constitutes the parameter vectors $\psi, \phi, \xi$.

The optimization is performed using the Adam algorithm \cite{kingma2014adam}, which is a form of stochastic gradient descent with adaptive momentum. Given the total loss function $\mathcal{L}(\paramvec)$, the parameter update at iteration $k$ follows: $\paramvec_{k+1} = \paramvec_k - \eta \hat{\bm{m}}_k / (\sqrt{\hat{\bm{v}}_k} + \epsilon)$, where $\eta$ is the learning rate (set to $10^{-4}$ in our experiments) and $\hat{\bm{m}}_k, \hat{\bm{v}}_k$ are bias-corrected estimates of the first and second moments of the gradient $\nabla_{\paramvec} \mathcal{L}(\paramvec_k)$.

\subsection{SDE Discretization and Monte Carlo Estimation}
The expectations in the loss functions are approximated via Monte Carlo integration over trajectories of the state variables. These trajectories are generated by discretizing the governing SDEs. Let the time horizon $[0, T]$ be divided into $N$ steps of size $\Delta t = T/N$, creating a time grid $t_k = k\Delta t$ for $k=0, \dots, N$. The continuous-time state process is
\begin{equation}
d\statevec_t = \driftvec_{\statevec}(t, \statevec_t, \controlvec_t) dt + \diffvec_{\statevec}(t, \statevec_t, \controlvec_t) d\bm{M}_t,
\end{equation}
where $\controlvec_t = (\Theta_\psi(t,\statevec_t), \mathcal{A}_\phi(t,\statevec_t))$ is the policy function, and $\bm{M}_t$ is a continuous vector-valued martingale with quadratic variation process $d\langle \bm{M} \rangle_t = \bm{Q}(t, \statevec_t) dt$.

We use the Euler-Maruyama scheme for discretization. The state at time $t_{k+1}$ is approximated as:
\begin{equation}
\statevec_{t_{k+1}} \approx \statevec_{t_k} + \driftvec_{\statevec}(t_k, \statevec_{t_k}, \controlvec_{t_k}) \Delta t + \diffvec_{\statevec}(t_k, \statevec_{t_k}, \controlvec_{t_k}) \Delta \bm{M}_{t_k}.
\end{equation}
The martingale increment $\Delta \bm{M}_{t_k} = \bm{M}_{t_{k+1}} - \bm{M}_{t_k}$ is a random variable sampled from a distribution with mean zero and covariance given by the discretized quadratic variation:
\begin{equation}
\E[\Delta \bm{M}_{t_k} \Delta \bm{M}_{t_k}^\top \mid \mathcal{F}_{t_k}] = \E \left[ \int_{t_k}^{t_{k+1}} \bm{Q}(s, \statevec_s) ds \mid \mathcal{F}_{t_k} \right] \approx \bm{Q}(t_k, \statevec_{t_k}) \Delta t.
\end{equation}
In practice, this means we sample $\Delta \bm{M}_{t_k}$ from a multivariate normal distribution $\mathcal{N}(\bm{0}, \bm{Q}(t_k, \statevec_{t_k}) \Delta t)$. The integrals in the loss functions are then approximated by sums over the discretized trajectory, and the expectation is approximated by the sample mean over a batch of $M$ such simulated trajectories (we use $M=1024$).

\subsection{Computation of Derivatives via Automatic Differentiation}
A critical component of the NBO method is the computation of the value function's derivatives, $\partial_t V_\xi, \nabla_{\statevec} V_\xi, \nabla_{\statevec\statevec}^2 V_\xi$, which are required for the HJB loss. These are computed using automatic differentiation (AD), a set of techniques for numerically evaluating the derivative of a function specified by a computer program. AD leverages the fact that any program executes a sequence of elementary arithmetic operations. By applying the chain rule repeatedly to these operations, one can compute derivatives automatically and to machine precision.

The gradients $\partial_t V_\xi$ and $\nabla_{\statevec} V_\xi$ are computed efficiently via a single pass of the backpropagation algorithm. The Hessian matrix $\nabla_{\statevec\statevec}^2 V_\xi$ is more computationally intensive. Modern deep learning libraries (like PyTorch or TensorFlow) compute the term $\text{Tr}(\bm{\sigma} \bm{Q} \bm{\sigma}^\top \nabla_{\statevec\statevec}^2 V)$ efficiently, often without explicitly constructing the full Hessian matrix. This is typically achieved using Hessian-vector products, which scales more favorably with the state dimension than forming the full Hessian.

\section{Hyperparameter Selection and Sensitivity}
\label{app:hyperparams}

\subsection{Theoretical Role and Sensitivity of the HJB Loss Weight (\texorpdfstring{$\omega$}{omega})}
The hyperparameter $\omega$ in the total loss function, $\mathcal{L} = L_{\text{PERF}} + \omega L_{\text{HJB}}$, is central to the actor-critic dynamic. To understand its role formally, consider the gradients with respect to the actor parameters $\bm{\pi} = (\psi, \phi)$ and critic parameters $\xi$:
\begin{align}
\nabla_{\bm{\pi}} \mathcal{L} &= \nabla_{\bm{\pi}} L_{\text{PERF}}(\bm{\pi}, \xi) + \omega \nabla_{\bm{\pi}} L_{\text{HJB}}(\bm{\pi}, \xi) \\
\nabla_\xi \mathcal{L} &= \nabla_\xi L_{\text{PERF}}(\bm{\pi}, \xi) + \omega \nabla_\xi L_{\text{HJB}}(\bm{\pi}, \xi)
\end{align}
The term $L_{\text{PERF}}$ is the negative of the expected Hamiltonian. In an actor-critic algorithm, the critic's parameters $\xi$ are primarily updated based on a measure of Bellman error, while the actor's parameters $\bm{\pi}$ are updated based on the performance gradient provided by the critic.

In our framework, the dominant term for the critic update is $\omega \nabla_\xi L_{\text{HJB}}$, as it directly penalizes deviations from the Bellman equation. The dominant term for the actor update is $\nabla_{\bm{\pi}} L_{\text{PERF}}$, which performs gradient ascent on the Hamiltonian using the critic's valuation. The parameter $\omega$ controls the relative scale of the gradients driving the critic's learning. A larger $\omega$ can be interpreted as increasing the critic's learning rate relative to the actor's.

This connects directly to the two-timescale stochastic approximation theory discussed in Appendix \ref{app:convergence}. For stable convergence of an actor-critic system, the critic should learn "faster" than the actor, allowing it to provide a stable evaluation of a given policy before the policy itself changes significantly. Setting $\omega > 1$ is a simple, direct mechanism to achieve this timescale separation. Our empirical choice of $\omega=1.0$ proved to be a good default, providing a natural normalization of the gradient magnitudes for the problems studied. To validate this choice, Figure \ref{fig:omega_sensitivity} shows a sensitivity analysis. While the convergence speed is affected by $\omega$, the algorithm is stable and converges to the correct solution for a range of values. The choice of $\omega=1.0$ represents a robust and well-balanced baseline.

\begin{figure}[H]
    \centering
    \begin{tikzpicture}
        \begin{axis}[
            title={Convergence Error vs. HJB Loss Weight $\omega$},
            xlabel={Training Epochs},
            ylabel={Mean Squared Error},
            ymode=log,
            legend pos=north east,
            width=10cm,
            height=7cm,
            ]
            \addplot[blue, thick, mark=*] table[x=epoch, y=error_w1] {omega_sensitivity.dat};
            \addlegendentry{$\omega = 1.0$ (Baseline)}
            \addplot[red, thick, mark=square*] table[x=epoch, y=error_w01] {omega_sensitivity.dat};
            \addlegendentry{$\omega = 0.1$}
            \addplot[green!50!black, thick, mark=triangle*] table[x=epoch, y=error_w10] {omega_sensitivity.dat};
            \addlegendentry{$\omega = 10.0$}
        \end{axis}
    \end{tikzpicture}
    \caption{Sensitivity of convergence to the HJB loss weight $\omega$ for the Merton problem. The y-axis plots a measure of solution error (mean squared error of policy function from analytical solution) on a log scale. While a very low weight ($\omega=0.1$) slows convergence, the algorithm is robust for $\omega=1.0$ and $\omega=10.0$, supporting the choice of $\omega=1.0$ as a reasonable default.}
    \label{fig:omega_sensitivity}
\end{figure}

\section{Formal Proofs}
\label{app:proofs}

This appendix provides a rigorous proof of Proposition \ref{prop:optimality}. We prove that if the NBO algorithm finds a global minimizer of the population loss function, the resulting value and policy networks correspond to the unique viscosity solution of the underlying stochastic control problem.

\subsection{Formal Setup and Assumptions}
We first restate the problem and formalize the assumptions required for the proof. The agent's control problem is characterized by the Hamilton-Jacobi-Bellman (HJB) equation:
\begin{equation} \label{eq:app_hjb_full_repeat}
-\partial_t V(t,\statevec) - \hamiltonian(t, \statevec, V, \nabla_{\statevec} V, \nabla_{\statevec\statevec}^2 V) = 0,
\end{equation}
with a terminal condition $V(T, \statevec) = G(\statevec)$. The Hamiltonian, $\hamiltonian$, is defined as the supremum of the pre-Hamiltonian, $\mathcal{H}^{\controlvec}$, over the compact set of feasible controls $\mathcal{K}$:
\begin{equation}
\hamiltonian(t, \statevec, v, \bm{p}, \bm{P}) := \sup_{\controlvec \in \mathcal{K}} \left\{ \mathcal{H}^{\controlvec}(t, \statevec, v, \bm{p}, \bm{P}) \right\},
\end{equation}
where the pre-Hamiltonian for a specific control $\controlvec$ is given by:
\begin{equation}
\mathcal{H}^{\controlvec}(t, \statevec, v, \bm{p}, \bm{P}) := f(t,\statevec,\controlvec,v) + \bm{p}^\top \driftvec_{\statevec}(t,\statevec,\controlvec) + \frac{1}{2}\text{Tr}\left(\diffvec_{\statevec}(t,\statevec,\controlvec) \bm{Q}_t \diffvec_{\statevec}(t,\statevec,\controlvec)^\top \bm{P}\right).
\end{equation}
Here, $v$ is a scalar representing the value, $\bm{p}$ is a vector representing its gradient, and $\bm{P}$ is a matrix representing its Hessian.

We proceed under the following standard assumptions.

\begin{assumption}[Regularity of Problem Data]
\label{ass:regularity_app}
The functions defining the problem—the utility aggregator $f$, the drift $\driftvec_{\statevec}$, the diffusion $\diffvec_{\statevec}$, the quadratic variation matrix $\bm{Q}_t$, and the terminal value function $G$—are continuous in all their arguments. Furthermore, they are Lipschitz continuous with respect to the state $\statevec$, control $\controlvec$, and value $v$ arguments. The control set $\mathcal{K}$ is compact.
\end{assumption}

\begin{remark}
Under Assumption \ref{ass:regularity_app}, the existence of a unique continuous viscosity solution $V^*(t,\statevec)$ to the HJB equation \eqref{eq:app_hjb_full_repeat} is guaranteed (see, e.g., \cite{fleming2006controlled}, Theorem IV.3.1 and related results).
\end{remark}

\begin{assumption}[Universal Approximation of Functions and Derivatives]
\label{ass:uat_app}
The function classes for the value network, $\mathcal{F}_\xi$, and the policy networks, $\mathcal{F}_{\bm{\pi}}$, are universal approximators for functions and their derivatives on compact domains. Specifically, for any target function in $C^2$, there exists a network in the class that can approximate the function and its derivatives up to second order arbitrarily well. This requires the use of sufficiently smooth network activation functions (e.g., of class $C^2$) as formally established in \cite{hornik1990universal}. Our implementation adheres to this by using a $C^\infty$ activation function (GELU).
\end{assumption}

\begin{assumption}[Global Minimizer]
\label{ass:global_min_app}
Let $(\bm{\pi}^*, \xi^*)$ be a set of parameters corresponding to a global minimum of the population loss function $\mathcal{L}(\bm{\pi}, \xi) = L_{\text{PERF}}(\bm{\pi}, \xi) + \omega L_{\text{HJB}}(\bm{\pi}, \xi)$. We analyze the properties of the corresponding functions $V_{\xi^*}$ and $\controlvec^{\bm{\pi}^*}$.
\end{assumption}

\subsection{Proof of Proposition 1}

The proof proceeds by demonstrating that the conditions implied by a global minimum of the NBO loss function are precisely the conditions that define the unique viscosity solution of the HJB equation.

\paragraph{Step 1: The Critic's Equilibrium Condition.}
The total loss function $\mathcal{L}$ is a sum of the performance loss and the HJB loss, weighted by $\omega > 0$. The HJB loss, defined as an expectation of a squared residual, is non-negative: $L_{\text{HJB}} \ge 0$. A global minimum of $\mathcal{L}$ requires that each non-negative component contributing to it must be at its minimum possible value. Therefore, at the global optimum $(\bm{\pi}^*, \xi^*)$, the HJB loss must be zero:
\begin{equation}
L_{\text{HJB}}(\bm{\pi}^*, \xi^*) = \E \left[ \left( -\partial_t V_{\xi^*} - \mathcal{H}^{\bm{\pi}^*}(t, \statevec, V_{\xi^*}, \nabla V_{\xi^*}, \nabla^2 V_{\xi^*}) \right)^2 \right] = 0.
\end{equation}
The expectation is taken over the distribution of state-time points $(t, \statevec)$ sampled by the algorithm. For the expectation of a non-negative quantity to be zero, the quantity itself must be zero almost everywhere with respect to the sampling measure. By the continuity of the neural networks and the problem data (Assumptions \ref{ass:regularity_app} and \ref{ass:uat_app}), this implies that the learned value function $V_{\xi^*}$ satisfies the following partial differential equation pointwise for the learned policy $\controlvec^{\bm{\pi}^*}$:
\begin{equation} \label{eq:app_critic_eq_final}
-\partial_t V_{\xi^*}(t,\statevec) = \mathcal{H}^{\controlvec^{\bm{\pi}^*}}(t, \statevec, V_{\xi^*}, \nabla_{\statevec} V_{\xi^*}, \nabla_{\statevec\statevec}^2 V_{\xi^*}).
\end{equation}
This condition states that the critic network $V_{\xi^*}$ has learned to be a perfect Bellman-consistent valuation of the actor's policy $\controlvec^{\bm{\pi}^*}$.

\paragraph{Step 2: The Actor's Equilibrium Condition.}
Simultaneously, the parameters $(\bm{\pi}^*, \xi^*)$ must minimize $\mathcal{L}$ with respect to the actor parameters $\bm{\pi}$. This is equivalent to maximizing the performance objective, which, by the principle of optimality, means that the policy $\controlvec^{\bm{\pi}^*}$ must be the optimal policy given the valuation provided by the critic $V_{\xi^*}$. If it were not, there would exist another feasible policy $\tilde{\controlvec}$ and a set of states with positive measure where using $\tilde{\controlvec}$ would increase the Hamiltonian's value. This would lead to a higher total lifetime utility (and lower performance loss), contradicting the assumption that $(\bm{\pi}^*, \xi^*)$ is a global minimum.
Therefore, the learned policy must maximize the pre-Hamiltonian pointwise for all $(t, \statevec)$:
\begin{equation}
\controlvec^{\bm{\pi}^*}(t, \statevec) \in \argmax_{\controlvec \in \mathcal{K}} \left\{ \mathcal{H}^{\controlvec}(t, \statevec, V_{\xi^*}, \nabla_{\statevec} V_{\xi^*}, \nabla_{\statevec\statevec}^2 V_{\xi^*}) \right\}.
\end{equation}
This implies that the pre-Hamiltonian evaluated at the learned policy is equal to the full, supremized Hamiltonian:
\begin{equation} \label{eq:app_actor_eq_final}
\mathcal{H}^{\controlvec^{\bm{\pi}^*}}(t, \statevec, V_{\xi^*}, \dots) = \sup_{\controlvec \in \mathcal{K}} \mathcal{H}^{\controlvec}(t, \statevec, V_{\xi^*}, \dots) = \hamiltonian(t, \statevec, V_{\xi^*}, \dots).
\end{equation}

\paragraph{Step 3: Verification of the Viscosity Solution Property.}
We now combine these two equilibrium conditions to demonstrate that $V_{\xi^*}$, which is continuous by construction (as the output of a neural network with continuous activations), is indeed the unique viscosity solution of the HJB equation \eqref{eq:app_hjb_full_repeat}.

\begin{enumerate}[label=(\roman*)]
    \item \textbf{Subsolution property.} We must show that for any smooth test function $\varphi \in C^{1,2}$ such that $V_{\xi^*} - \varphi$ attains a local maximum at a point $(t_0, \statevec_0)$, the following inequality holds:
    \begin{equation*}
        -\partial_t \varphi(t_0, \statevec_0) - \hamiltonian(t_0, \statevec_0, V_{\xi^*}(t_0, \statevec_0), \nabla_{\statevec} \varphi(t_0, \statevec_0), \nabla_{\statevec\statevec}^2 \varphi(t_0, \statevec_0)) \le 0.
    \end{equation*}
    At the local maximum $(t_0, \statevec_0)$, we have $V_{\xi^*}(t_0, \statevec_0) = \varphi(t_0, \statevec_0)$ and the derivatives satisfy $\partial_t V_{\xi^*} \approx \partial_t \varphi$, $\nabla_{\statevec} V_{\xi^*} \approx \nabla_{\statevec} \varphi$, and $\nabla_{\statevec\statevec}^2 V_{\xi^*} \le \nabla_{\statevec\statevec}^2 \varphi$ in the matrix sense.
    
    From the critic's equilibrium condition \eqref{eq:app_critic_eq_final} at $(t_0, \statevec_0)$, we can write:
    \begin{align}
        -\partial_t \varphi(t_0, \statevec_0) &\approx -\partial_t V_{\xi^*}(t_0, \statevec_0) \nonumber \\
        &= \mathcal{H}^{\controlvec^{\bm{\pi}^*}}(t_0, \statevec_0, V_{\xi^*}, \nabla_{\statevec} V_{\xi^*}, \nabla_{\statevec\statevec}^2 V_{\xi^*}) \nonumber \\
        &\approx \mathcal{H}^{\controlvec^{\bm{\pi}^*}}(t_0, \statevec_0, \varphi, \nabla_{\statevec} \varphi, \nabla_{\statevec\statevec}^2 V_{\xi^*}). \label{eq:subsol_step1}
    \end{align}
    The Hamiltonian is monotone increasing in its Hessian argument because the term $\frac{1}{2}\text{Tr}(\diffvec \bm{Q} \diffvec^\top \bm{P})$ involves the trace of a product of positive semi-definite matrices. Since $\bm{Q}_t$ and $\diffvec \diffvec^\top$ are positive semi-definite, their product is also positive semi-definite, and the trace operator is a linear functional with positive semi-definite arguments. Since $\nabla_{\statevec\statevec}^2 V_{\xi^*} \le \nabla_{\statevec\statevec}^2 \varphi$ in the matrix sense, we have:
    \begin{equation}
        \mathcal{H}^{\controlvec^{\bm{\pi}^*}}(t_0, \statevec_0, \varphi, \nabla_{\statevec} \varphi, \nabla_{\statevec\statevec}^2 V_{\xi^*}) \le \mathcal{H}^{\controlvec^{\bm{\pi}^*}}(t_0, \statevec_0, \varphi, \nabla_{\statevec} \varphi, \nabla_{\statevec\statevec}^2 \varphi). \label{eq:subsol_step2}
    \end{equation}
    Furthermore, by definition of the supremum, the pre-Hamiltonian for any policy is less than or equal to the full Hamiltonian:
    \begin{equation}
        \mathcal{H}^{\controlvec^{\bm{\pi}^*}}(t_0, \statevec_0, \varphi, \nabla_{\statevec} \varphi, \nabla_{\statevec\statevec}^2 \varphi) \le \hamiltonian(t_0, \statevec_0, \varphi, \nabla_{\statevec} \varphi, \nabla_{\statevec\statevec}^2 \varphi). \label{eq:subsol_step3}
    \end{equation}
    Combining equations \eqref{eq:subsol_step1}, \eqref{eq:subsol_step2}, and \eqref{eq:subsol_step3} gives the desired result:
    \begin{equation*}
         -\partial_t \varphi(t_0, \statevec_0) \le \hamiltonian(t_0, \statevec_0, \varphi, \nabla_{\statevec} \varphi, \nabla_{\statevec\statevec}^2 \varphi),
    \end{equation*}
    which is the subsolution property.

    \item \textbf{Supersolution property.} We must show that for any smooth test function $\varphi \in C^{1,2}$ such that $V_{\xi^*} - \varphi$ has a local minimum at $(t_0, \statevec_0)$, the following inequality holds:
    \begin{equation*}
        -\partial_t \varphi(t_0, \statevec_0) - \hamiltonian(t_0, \statevec_0, V_{\xi^*}(t_0, \statevec_0), \nabla_{\statevec} \varphi(t_0, \statevec_0), \nabla_{\statevec\statevec}^2 \varphi(t_0, \statevec_0)) \ge 0.
    \end{equation*}
    At the local minimum $(t_0, \statevec_0)$, we have $V_{\xi^*}(t_0, \statevec_0) = \varphi(t_0, \statevec_0)$ and the derivatives satisfy $\partial_t V_{\xi^*} \approx \partial_t \varphi$, $\nabla_{\statevec} V_{\xi^*} \approx \nabla_{\statevec} \varphi$, and $\nabla_{\statevec\statevec}^2 V_{\xi^*} \ge \nabla_{\statevec\statevec}^2 \varphi$.
    
    We again start from the critic's equilibrium condition \eqref{eq:app_critic_eq_final}, but now combine it with the actor's equilibrium condition \eqref{eq:app_actor_eq_final}. Crucially, this step relies on both conditions holding simultaneously to fully enforce the supremum operator present in the HJB equation.
    \begin{align}
        -\partial_t \varphi(t_0, \statevec_0) &\approx -\partial_t V_{\xi^*}(t_0, \statevec_0) \nonumber \\
        &= \mathcal{H}^{\controlvec^{\bm{\pi}^*}}(t_0, \statevec_0, V_{\xi^*}, \nabla_{\statevec} V_{\xi^*}, \nabla_{\statevec\statevec}^2 V_{\xi^*}) \nonumber \\
        &= \hamiltonian(t_0, \statevec_0, V_{\xi^*}, \nabla_{\statevec} V_{\xi^*}, \nabla_{\statevec\statevec}^2 V_{\xi^*}). \label{eq:supersol_step1}
    \end{align}
    Substituting the properties at the minimum point, we have:
    \begin{equation}
         -\partial_t \varphi(t_0, \statevec_0) \approx \hamiltonian(t_0, \statevec_0, \varphi, \nabla_{\statevec} \varphi, \nabla_{\statevec\statevec}^2 V_{\xi^*}). \label{eq:supersol_step2}
    \end{equation}
    Using the monotonicity of the Hamiltonian in its Hessian argument and the fact that $\nabla_{\statevec\statevec}^2 V_{\xi^*} \ge \nabla_{\statevec\statevec}^2 \varphi$, we obtain:
    \begin{equation}
        \hamiltonian(t_0, \statevec_0, \varphi, \nabla_{\statevec} \varphi, \nabla_{\statevec\statevec}^2 V_{\xi^*}) \ge \hamiltonian(t_0, \statevec_0, \varphi, \nabla_{\statevec} \varphi, \nabla_{\statevec\statevec}^2 \varphi). \label{eq:supersol_step3}
    \end{equation}
    Combining \eqref{eq:supersol_step2} and \eqref{eq:supersol_step3} gives the supersolution property:
    \begin{equation*}
        -\partial_t \varphi(t_0, \statevec_0) \ge \hamiltonian(t_0, \statevec_0, \varphi, \nabla_{\statevec} \varphi, \nabla_{\statevec\statevec}^2 \varphi).
    \end{equation*}
\end{enumerate}

\paragraph{Step 4: Uniqueness and Conclusion.}
The function $V_{\xi^*}$ is continuous and has been shown to satisfy both the subsolution and supersolution properties at any point in its domain. By definition, $V_{\xi^*}$ is a viscosity solution of the HJB equation \eqref{eq:app_hjb_full_repeat}.

Under Assumption \ref{ass:regularity_app}, the viscosity solution to this HJB equation is unique. Therefore, it must be that the value function learned by the NBO algorithm, $V_{\xi^*}$, is identical to the true value function of the control problem, $V^*$. Consequently, the policy $\controlvec^{\bm{\pi}^*}$, which by construction is the optimal policy for the value function $V_{\xi^*} = V^*$, is the true optimal policy.

This completes the proof that a global minimizer of the NBO loss function corresponds to the unique, correct (viscosity) solution of the stochastic optimal control problem.


\section{Theoretical Underpinnings of the NBO Algorithm}
\label{app:uat_conv}

This section provides the formal theoretical basis for the NBO method. We first establish that neural networks possess the necessary approximation capabilities for both the value function and its derivatives, a prerequisite for the validity of the HJB-based loss. We then analyze the convergence properties of the actor-critic optimization algorithm using the framework of two-timescale stochastic approximation.

\subsection{Approximation Capabilities of Neural Networks}
\label{app:uat}

The validity of minimizing the HJB residual via the loss function $L_{\text{HJB}}$ hinges on the capacity of neural networks to approximate not just the value function itself, but also its partial derivatives up to second order. This is a stronger condition than that provided by classical universal approximation theorems.

The foundational result, the Universal Approximation Theorem (UAT) by \cite{hornik1989multilayer}, states that a standard feedforward network with a single hidden layer can approximate any continuous function on a compact domain. However, this does not guarantee the approximation of derivatives. A more powerful result is required.

\begin{theorem}[Universal Approximation of Functions and Their Derivatives, \cite{hornik1990universal,qian2025icml}]
\label{thm:uat_derivatives}
Let $\sigma(\cdot)$ be a non-polynomial activation function that is continuously differentiable up to order $m \ge 1$ (i.e., $\sigma \in C^m(\mathbb{R})$). Let $\mathcal{F}_{d,m}$ be the class of functions $f: \mathbb{R}^d \to \mathbb{R}$ that are in $C^m$ and for which the function itself and all its partial derivatives up to order $m$ are bounded and continuous. Then, for any target function $f \in \mathcal{F}_{d,m}$, any compact set $K \subset \mathbb{R}^d$, and any $\epsilon > 0$, there exists a feedforward neural network $\hat{f}$ with one hidden layer using activation function $\sigma$ such that:
\begin{equation}
\sup_{x \in K} \left| D^\alpha f(x) - D^\alpha \hat{f}(x) \right| < \epsilon,
\end{equation}
for all multi-indices $\alpha = (\alpha_1, \dots, \alpha_d)$ with $|\alpha| = \sum_{i=1}^d \alpha_i \le m$.
\end{theorem}

\subsubsection{Implications for the NBO Method}
The HJB equation \eqref{eq:hjb} is a second-order non-linear partial differential equation. Its residual, which is the target of the critic's loss $L_{\text{HJB}}$, explicitly involves the value function $V_\xi$, its gradient $\nabla_{\statevec} V_\xi$, and its Hessian matrix $\nabla_{\statevec\statevec}^2 V_\xi$. To ensure that a zero HJB loss corresponds to a true solution, the network $V_\xi$ must be able to simultaneously approximate the true value function $V^*$ (zeroth order), its gradient $\nabla V^*$ (first order), and its Hessian $\nabla^2 V^*$ (second order).

Theorem \ref{thm:uat_derivatives} provides this guarantee. By choosing a network architecture with an activation function of class $C^2$ (e.g., the hyperbolic tangent, \textit{tanh}, or Gaussian Error Linear Unit, GELU), the theorem ensures that a sufficiently large network can indeed make the approximation error for the function and its first two derivatives arbitrarily small on any compact domain. This guarantees that the loss landscape has a global minimum at zero, which corresponds to a function that genuinely solves the PDE, rather than a spurious solution where the residual is minimized due to fortuitous cancellation of poorly approximated derivative terms.

\begin{remark}[Practical Choice of Activation Function]
\label{rem:relu_practical}
As stated in Section \ref{sec:theoretical_framework} and implemented in Appendix \ref{app:implementation}, our work uses smooth ($C^\infty$) GELU activation functions to maintain full theoretical consistency. Historically and in many applied deep learning domains, the non-smooth Rectified Linear Unit (ReLU), $g(x)=\max(0,x)$, is popular for its computational speed. While empirical work has shown ReLU-based networks can be surprisingly effective at solving PDEs, this success relies on implicit regularization properties of stochastic optimization rather than on the formal approximation guarantees of Theorem \ref{thm:uat_derivatives}. For the level of rigor demanded by the economic problems we address, choosing an activation function that satisfies the premises of the underlying theory is the appropriate and necessary choice.
\end{remark}

\subsection{Convergence Analysis via Two-Timescale Stochastic Approximation}
\label{app:convergence}

Proposition \ref{prop:optimality} establishes that the NBO loss function is a correct objective whose global minimizer is the solution to the economic problem. We now provide a formal argument for why the optimization algorithm itself—simultaneous stochastic gradient descent on the actor and critic parameters—can be expected to converge to a meaningful solution. The appropriate mathematical framework is that of two-timescale stochastic approximation (TTSA), as developed in works such as \cite{borkar2008stochastic}.

\subsubsection{Stochastic Approximation Formulation}
Let $\paramvec_k = (\bm{\pi}_k, \xi_k)$ represent the concatenated actor and critic parameters at iteration $k$. The stochastic gradient descent (SGD) update rule with learning rates $\eta_{\bm{\pi}}(k)$ and $\eta_\xi(k)$ is:
\begin{align}
\bm{\pi}_{k+1} &= \bm{\pi}_k - \eta_{\bm{\pi}}(k) \left( \nabla_{\bm{\pi}} \mathcal{L}(\bm{\pi}_k, \xi_k) + \bm{M}_{\pi, k+1} \right) \\
\xi_{k+1} &= \xi_k - \eta_\xi(k) \left( \nabla_\xi \mathcal{L}(\bm{\pi}_k, \xi_k) + M_{\xi, k+1} \right)
\end{align}
where $\mathcal{L}$ is the empirical loss on a randomly sampled mini-batch, and $\bm{M}_{\pi, k+1}, M_{\xi, k+1}$ are zero-mean martingale difference noise terms arising from the Monte Carlo sampling. For convergence, the learning rates must satisfy the standard Robbins-Monro conditions: $\sum_{k=0}^\infty \eta(k) = \infty$ and $\sum_{k=0}^\infty \eta(k)^2 < \infty$.

\subsubsection{The Two-Timescale Argument for Actor-Critic Stability}
The stability of the actor-critic learning process depends on a separation of learning speeds. The critic should update "faster" than the actor, allowing it to converge to an accurate evaluation of the actor's current policy before that policy changes significantly. This prevents the actor from chasing a noisy, rapidly moving target. In our framework, this can be achieved by setting the learning rates such that $\eta_{\bm{\pi}}(k) = o(\eta_\xi(k))$, or implicitly by setting the HJB loss weight $\omega \gg 1$.

The theory of stochastic approximation asserts that the asymptotic behavior of such a discrete-time iterative process is described by a coupled system of Ordinary Differential Equations (ODEs). Due to the timescale separation, this system can be analyzed by considering the fast and slow dynamics separately.

\paragraph{Fast Timescale (The Critic's Dynamics).}
For a "frozen" actor policy with parameters $\bm{\pi}$, the critic's parameters $\xi$ evolve on a fast timescale $\tau$. Their dynamics are governed by the ODE:
\begin{equation}
\frac{d\xi(\tau)}{d\tau} = -\nabla_\xi \mathbb{E}[\mathcal{L}(\bm{\pi}, \xi(\tau))] \approx -\omega \nabla_\xi \mathbb{E}[L_{\text{HJB}}(\bm{\pi}, \xi(\tau))].
\end{equation}
The approximation holds because the critic's parameters $\xi$ primarily affect the HJB loss. Under standard regularity conditions, this gradient flow dynamic converges to an equilibrium point (or set) $\bar{\xi}(\bm{\pi})$ where the gradient is zero. This equilibrium is precisely where $\nabla_\xi L_{\text{HJB}}(\bm{\pi}, \bar{\xi}(\bm{\pi}))=0$, meaning the critic has converged to the true, Bellman-consistent value function for the fixed policy specified by $\bm{\pi}$.

\paragraph{Slow Timescale (The Actor's Dynamics).}
The actor's parameters $\bm{\pi}$ evolve on a slower timescale $t$. The TTSA framework allows us to analyze their dynamics by assuming the critic is always at its fast-timescale equilibrium, i.e., $\xi_k \approx \bar{\xi}(\bm{\pi}_k)$. The actor's ODE is then:
\begin{equation}
\frac{d\bm{\pi}(t)}{dt} = -\nabla_{\bm{\pi}} \mathbb{E}[\mathcal{L}(\bm{\pi}(t), \bar{\xi}(\bm{\pi}(t)))].
\end{equation}
At the critic's equilibrium, $L_{\text{HJB}}(\bm{\pi}, \bar{\xi}(\bm{\pi})) \approx 0$. Therefore, the total loss is dominated by the performance loss: $\mathcal{L}(\bm{\pi}, \bar{\xi}(\bm{\pi})) \approx L_{\text{PERF}}(\bm{\pi}, \bar{\xi}(\bm{\pi}))$. Let $V^{\bm{\pi}}(\cdot) \equiv V_{\bar{\xi}(\bm{\pi})}(\cdot)$ denote the true value function for policy $\bm{\pi}$. By definition, the performance loss is the negative of the expected Hamiltonian, whose integral is the total value. The actor's ODE thus simplifies to a policy gradient update:
\begin{equation}
\frac{d\bm{\pi}(t)}{dt} \approx - \nabla_{\bm{\pi}} \left( -V^{\bm{\pi}}(0, \statevec_0) \right) = \nabla_{\bm{\pi}} V^{\bm{\pi}}(0, \statevec_0).
\end{equation}
The actor continuously adjusts its policy in the direction of steepest ascent of the true value function.

\subsubsection{Convergence to a Nash Equilibrium}
The TTSA theory guarantees that the joint process $(\bm{\pi}_k, \xi_k)$ converges (in a weak sense) to a stationary point $(\bm{\pi}^*, \xi^*)$ of the coupled ODE system. A stationary point is one where both dynamics are at a fixed point: the critic is at equilibrium for the actor's policy, and the actor's policy is at a local optimum given the critic's valuation. This stationary point corresponds to a local Nash Equilibrium of the game between the actor and the critic. While this analysis does not guarantee convergence to the global optimum due to the non-convexity of the loss landscape, it provides a rigorous theoretical foundation for the algorithm's observed stability and its convergence to a locally optimal and internally consistent solution.


\section{Robustness and Reproducibility}
\label{app:robustness}
A crucial aspect of establishing the practical utility of any numerical method is demonstrating its robustness to sources of stochasticity inherent in the algorithm. For the NBO method, the primary source of randomness is the initialization of the neural network weights and the stochastic nature of the Monte Carlo sampling used in the loss function.

To validate the stability and reliability of our results, all benchmark experiments presented in Section \ref{sec:validation} and \ref{sec:recursive_test} were performed multiple times. Specifically, as noted in Tables \ref{tab:merton_results}, \ref{tab:merton_constrained}, and \ref{tab:ez_results}, each validation exercise was conducted over 10 independent training runs, each starting from a different random seed. This ensures that our results are not an artifact of a single "lucky" weight initialization. The reported means and standard deviations of the learned policies show that the algorithm consistently converges to the same, correct solution with very low variance across runs. This provides strong empirical evidence that the local Nash Equilibria found by the stochastic gradient descent optimizer correspond to the true global optimum for these problems, supporting the claims made in Section \ref{sec:theoretical_framework}. The full code provided for replication allows for these robustness checks to be verified and extended.

\end{appendix}



\phantomsection
\addcontentsline{toc}{section}{References}
\bibliography{reference}
\bibliographystyle{ecta-fullname} 


\clearpage
\section*{Supplementary Appendix to “Neural Bellman Operators” (Not For Publication)}\label{sec:olappendix} 
\setcounter{equation}{0} 
\renewcommand{\theequation}{S\arabic{equation}} 
\setcounter{assumption}{0} 
\renewcommand{\theassumption}{S\arabic{assumption}} 
\setcounter{proposition}{0} 
\renewcommand{\theproposition}{S\arabic{proposition}} 
\setcounter{remark}{0} 
\renewcommand{\theremark}{S\arabic{remark}} 
\setcounter{table}{0} 
\renewcommand{\thetable}{S\arabic{table}} 
\setcounter{figure}{0} 
\renewcommand{\thefigure}{S\arabic{figure}} 
\setcounter{section}{0} 
\renewcommand{\thesection}{S\arabic{section}} 


\section{Introduction}

This supplementary document provides a detailed formalization of the extensions of the Neural Bellman Operators (NBO) method discussed in Section 8 of the main paper, "Neural Bellman Operators." We present the precise mathematical formulation and the detailed algorithmic implementation for solving two classes of models that are central to modern economics but computationally challenging: (1) problems with time-inconsistent preferences, focusing on sophisticated quasi-hyperbolic agents, and (2) high-dimensional, continuous-time dynamic games, focusing on the computation of Markov Perfect Nash Equilibria (MPNE).

For each case, we first define the economic problem and its equilibrium concept, characterized by a generalized or coupled system of Hamilton-Jacobi-Bellman (HJB) equations. We then present a step-by-step NBO algorithm designed to find the specific equilibrium structure of that problem. This includes the parameterization of policies and value functions with neural networks, the definition of the specialized loss functions, and the logic of the optimization procedure. These sections provide the necessary detail for researchers wishing to implement these extensions.

\section{A Formal Framework for Time-Inconsistent Agents}
\label{sec:supp_time_inconsistency}

We formalize the NBO framework for finding the subgame-perfect equilibrium consumption plan for a sophisticated agent with quasi-hyperbolic preferences in continuous time. Such an agent understands their future selves will also be present-biased and acts accordingly. The equilibrium is an intra-personal Nash equilibrium between the agent's successive temporal selves.

\subsection{The Model and the Generalized HJB Equation}

Consider an agent with a time-inconsistent objective. The behavior of a sophisticated agent is characterized by a time-consistent policy $c^*(t, \statevec_t)$ that constitutes a subgame-perfect equilibrium. This equilibrium policy and its associated value function $V(t, \statevec_t)$ must satisfy a \textit{generalized} HJB equation (\cite{harris2003stochastic}).

Let the state vector $\statevec_t \in \R^d$ evolve according to the SDE:
\begin{equation}
d\statevec_t = \bm{\mu}(t, \statevec_t, c_t) dt + \bm{\sigma}(t, \statevec_t, c_t) d\bm{M}_t,
\end{equation}
where $c_t$ is the control (e.g., consumption). The value function $V(t,\statevec_t)$ of a sophisticated agent is the lifetime utility evaluated from that point forward, *given* that all future selves will follow the equilibrium policy $c^*(s, \statevec_s)$ for $s > t$. The current self, however, gets to choose $c_t$ optimally, taking the future selves' actions as given.

\begin{definition}[Generalized HJB for Sophisticated Agents]
The value function $V(t,\statevec)$ and the equilibrium policy $c^*(t,\statevec)$ for a sophisticated agent solve the following fixed-point problem. The equilibrium condition is a fixed point: the policy that solves the maximization problem for the current self must be identical to the anticipated policy. At equilibrium, the value function solves:
\begin{equation} \label{eq:supp_ghjb_eq}
-\partial_t V(t,\statevec) = u(c^*) + \mathcal{D}^{c^*}V(t,\statevec),
\end{equation}
where the equilibrium policy $c^*$ is the solution to the maximization of the current self's perceived Hamiltonian. For a simple continuous-time analogue of $(\beta, \delta)$ discounting with instantaneous utility $u(c)$ and present-bias parameter $\beta$, the current self's optimization is:
\begin{equation} \label{eq:supp_ghjb_max}
c^*(t, \statevec) = \argmax_{c \in \mathcal{A}} \left\{ \beta u(c) + \nabla_{\statevec} V(t,\statevec)^\top \bm{\mu}(t,\statevec,c) \right\}.
\end{equation}
Note that the agent's current choice is based on a distorted instantaneous utility $\beta u(c)$, but the value function itself accumulates undiscounted instantaneous utility $u(c^*)$. This captures the essence of present bias where the agent acts impatiently now but evaluates their life from a time-consistent perspective.
\end{definition}

\subsection{The NBO Algorithm for Time-Inconsistency}

To solve this fixed-point problem, we propose a "two-policy" NBO algorithm. This algorithm uses three neural networks to find the equilibrium by explicitly modeling the acting self and the anticipated future selves.

\begin{enumerate}
    \item \textbf{Acting Policy Network $\mathcal{C}_{\phi}(t,\statevec)$:} Represents the current self's choice of consumption, $c_t$. Its parameters are $\phi$.
    \item \textbf{Future Policy Network $\mathcal{C}_{\hat{\phi}}(t,\statevec)$:} Represents the anticipated equilibrium policy of all future selves, $\hat{c}_t$. Its parameters are $\hat{\phi}$. This network acts as a "target" that evolves slowly.
    \item \textbf{Value Network $V_{\xi}(t,\statevec)$:} Approximates the equilibrium value function. Its parameters are $\xi$.
\end{enumerate}

The algorithm minimizes a composite loss function designed to drive the system to the equilibrium fixed point.

\begin{definition}[NBO Loss for Time-Inconsistency]
The total loss function is a weighted sum of three components:
\begin{equation}
\mathcal{L}(\phi, \hat{\phi}, \xi) = L_{\text{PERF}}(\phi, \xi) + \omega L_{\text{HJB}}(\hat{\phi}, \xi) + \lambda_{EQ} L_{\text{EQ}}(\phi, \hat{\phi}).
\end{equation}
The components are:
\begin{enumerate}
    \item \textbf{Performance Loss ($L_{\text{PERF}}$):} This is the actor loss for the acting policy network $\mathcal{C}_\phi$. It encourages the current self to maximize its perceived Hamiltonian from \eqref{eq:supp_ghjb_max}.
    \begin{equation}
    L_{\text{PERF}} = - \E \left[ \beta u(\mathcal{C}_{\phi}) + \nabla_{\statevec} V_\xi^\top \bm{\mu}(t,\statevec, \mathcal{C}_{\phi}) \right].
    \end{equation}
    The gradient $\nabla V_\xi$ is taken from the critic but treated as a fixed input for this optimization step.
    
    \item \textbf{Generalized HJB Loss ($L_{\text{HJB}}$):} This is the critic loss, which enforces the Bellman equation \eqref{eq:supp_ghjb_eq} for the future (target) policy.
    \begin{equation}
    L_{\text{HJB}} = \E \left[ \left( -\partial_t V_\xi - \left( u(\mathcal{C}_{\hat{\phi}}) + \mathcal{D}^{\mathcal{C}_{\hat{\phi}}}V_\xi \right) \right)^2 \right].
    \end{equation}
    This loss forces $V_\xi$ to be the correct value function for the policy $\mathcal{C}_{\hat{\phi}}$.
    
    \item \textbf{Equilibrium Consistency Loss ($L_{\text{EQ}}$):} This term enforces the fixed-point condition by penalizing differences between the acting policy and the future policy.
    \begin{equation}
    L_{\text{EQ}} = \E \left[ \left\| \mathcal{C}_{\phi}(t,\statevec_t) - \mathcal{C}_{\hat{\phi}}(t,\statevec_t) \right\|^2 \right].
    \end{equation}
\end{enumerate}
\end{definition}

The full algorithm is detailed in Algorithm \ref{alg:ti}. A key feature is the slow update of the future policy network's parameters $\hat{\phi}$ using Polyak averaging. This stabilizes training by preventing the equilibrium target from moving too quickly.

\begin{algorithm}[H]
\caption{NBO Algorithm for Sophisticated Time-Inconsistent Agents}
\label{alg:ti}
\begin{algorithmic}
\State \textbf{Initialize} acting policy network $\mathcal{C}_{\phi}$, value network $V_{\xi}$.
\State \textbf{Initialize} future policy network $\mathcal{C}_{\hat{\phi}}$ with same weights as $\mathcal{C}_{\phi}$ (i.e., $\hat{\phi} \leftarrow \phi$).
\State \textbf{Initialize} hyperparameters: learning rate $\eta$, loss weights $\omega, \lambda_{EQ}$, Polyak averaging rate $\tau \in (0,1)$.

\For{training iteration $k=1, 2, \dots$}
    \State Sample a batch of state-time points $\{(t_j, \statevec_j)\}_{j=1}^B$.
    
    \State \textbf{Compute Actions:} $c_j \leftarrow \mathcal{C}_{\phi}(t_j, \statevec_j)$, $\hat{c}_j \leftarrow \mathcal{C}_{\hat{\phi}}(t_j, \statevec_j)$
    
    \State \textbf{Compute Derivatives:} Compute $\partial_t V_\xi, \nabla_{\statevec} V_\xi, \nabla_{\statevec\statevec}^2 V_\xi$ via AD.
    
    \State \textbf{Compute Loss Components:}
    \State $L_{\text{PERF}} \leftarrow -\frac{1}{B}\sum_{j=1}^B \left( \beta u(c_j) + (\nabla_{\statevec} V_\xi(t_j,\statevec_j))^\top \bm{\mu}(t_j,\statevec_j, c_j) \right)$
    \State $L_{\text{HJB}} \leftarrow \frac{1}{B}\sum_{j=1}^B \left( -\partial_t V_\xi(t_j,\statevec_j) - \left( u(\hat{c}_j) + \mathcal{D}^{\hat{c}_j}V_\xi(t_j,\statevec_j) \right) \right)^2$
    \State $L_{\text{EQ}} \leftarrow \frac{1}{B}\sum_{j=1}^B \| c_j - \hat{c}_j \|^2$
    
    \State \textbf{Total Loss:} $\mathcal{L} \leftarrow L_{\text{PERF}} + \omega L_{\text{HJB}} + \lambda_{EQ} L_{\text{EQ}}$.
    
    \State \textbf{Update Acting and Critic Networks:}
    \State Update $\phi$ and $\xi$ using one step of gradient descent on $\mathcal{L}$:
    \State $(\phi, \xi) \leftarrow (\phi, \xi) - \eta \nabla_{(\phi, \xi)} \mathcal{L}$
    
    \State \textbf{Update Future Policy Network (Polyak Averaging):}
    \State $\hat{\phi} \leftarrow \tau \phi + (1-\tau) \hat{\phi}$
\EndFor
\State \textbf{Return} trained networks $(\mathcal{C}_{\phi}, V_{\xi})$, which approximate $(c^*, V)$.
\end{algorithmic}
\end{algorithm}


\section{A Formal Framework for Dynamic Games}
\label{sec:supp_games}

We now formalize the NBO framework for computing Markov Perfect Nash Equilibria (MPNE) in high-dimensional, continuous-time stochastic games with $N$ players.

\subsection{The Model and the System of Coupled HJB Equations}

Consider a game with $N$ players, indexed by $i \in \{1, \dots, N\}$. The state of the game is a vector $\statevec_t \in \R^d$ that may include player-specific components, $\statevec_t = (\statevec_t^1, \dots, \statevec_t^N)$. Player $i$ chooses an action $a_t^i \in \mathcal{A}^i$ to maximize their own objective functional, $V^i(t,\statevec_t)$. The evolution of the full state vector depends on the actions of all players, $\bm{a}_t = (a_t^1, \dots, a_t^N)$:
\begin{equation}
d\statevec_t = \bm{\mu}(t, \statevec_t, \bm{a}_t)dt + \bm{\sigma}(t, \statevec_t, \bm{a}_t)d\bm{M}_t.
\end{equation}

\begin{definition}[Markov Perfect Nash Equilibrium (MPNE)]
An MPNE is a set of policy functions (strategies) $\{\controlvec^{1,*}(t,\statevec), \dots, \controlvec^{N,*}(t,\statevec)\}$ such that for each player $i$, the policy $\controlvec^{i,*}$ is a best response to the other players' equilibrium policies $\{\controlvec^{j,*}\}_{j \neq i}$.
\end{definition}

The MPNE is characterized by a system of $N$ coupled HJB equations.

\begin{definition}[System of Coupled HJB Equations]
The set of value functions $\{V^1, \dots, V^N\}$ and policies $\{\controlvec^{1,*}, \dots, \controlvec^{N,*}\}$ at an MPNE must satisfy the following system of $N$ PDEs:
\begin{equation} \label{eq:supp_hjb_game}
-\partial_t V^i(t,\statevec) = \sup_{a^i \in \mathcal{A}^i} \left\{ f^i(\statevec, a^i, \controlvec^{-i,*}(\statevec)) + \mathcal{D}^{a^i, \controlvec^{-i,*}}V^i(t,\statevec) \right\},
\end{equation}
for each player $i=1, \dots, N$, with terminal conditions $V^i(T,\statevec) = G^i(\statevec)$. The It\^o operator for player $i$, $\mathcal{D}^{a^i, \controlvec^{-i,*}}$, is evaluated using player $i$'s choice $a^i$ and the other players' equilibrium policies $\controlvec^{-i,*}$.
\end{definition}

\subsection{The NBO Algorithm for Dynamic Games}

The NBO framework for games parameterizes the value function and policy function for each of the $N$ players with a dedicated pair of neural networks.

The framework requires $2N$ neural networks:
\begin{itemize}
    \item \textbf{Policy Networks:} $\{\mathcal{A}_{\phi_i}^i(t,\statevec)\}_{i=1}^N$.
    \item \textbf{Value Networks:} $\{V_{\xi_i}^i(t,\statevec)\}_{i=1}^N$.
\end{itemize}
The complete set of trainable parameters is $(\phi_1, \dots, \phi_N, \xi_1, \dots, \xi_N)$.

\begin{definition}[NBO Loss for Dynamic Games]
The global loss function is the sum of individual NBO loss functions for all players.
\begin{equation}
\mathcal{L}_{\text{Game}}(\{\phi_i, \xi_i\}_{i=1}^N) = \sum_{i=1}^N \mathcal{L}^i(\{\phi_j\}_{j=1}^N, \xi_i) = \sum_{i=1}^N \left( L_{\text{PERF}}^i + \omega_i L_{\text{HJB}}^i \right).
\end{equation}
For each player $i$, the components are:
\begin{enumerate}
    \item \textbf{Player $i$'s Performance Loss ($L_{\text{PERF}}^i$):} The actor loss for player $i$, encouraging policy $\mathcal{A}_{\phi_i}^i$ to maximize player $i$'s own Hamiltonian, evaluated by their own critic $V_{\xi_i}^i$.
    \begin{equation}
    L_{\text{PERF}}^i = - \E \left[ f^i(\statevec_t, \mathcal{A}_{\phi_i}^i, \{\mathcal{A}_{\phi_j}^j\}_{j \neq i}) + \mathcal{D}^{\mathcal{A}_{\phi_i}^i, \{\mathcal{A}_{\phi_j}^j\}_{j \neq i}}V^i_{\xi_i} \right].
    \end{equation}

    \item \textbf{Player $i$'s HJB Loss ($L_{\text{HJB}}^i$):} The critic loss for player $i$. It enforces player $i$'s HJB equation by matching the time derivative of the value function to the Hamiltonian evaluated at the current policies.
    \begin{equation}
    L_{\text{HJB}}^i = \E \left[ \left( -\partial_t V_{\xi_i}^i - \left( f^i(\cdot, \mathcal{A}_{\phi_i}^i, \{\mathcal{A}_{\phi_j}^j\}_{j\neq i}) + \mathcal{D}^{\mathcal{A}_{\phi_i}^i, \{\mathcal{A}_{\phi_j}^j\}_{j \neq i}}V_{\xi_i}^i \right) \right)^2 \right].
    \end{equation}
\end{enumerate}
\end{definition}

Minimizing this joint loss function drives the system to a fixed point where each player's policy is an approximate best response to the others. The full procedure is detailed in Algorithm \ref{alg:game}.

\begin{algorithm}[H]
\caption{NBO Algorithm for N-Player Dynamic Games}
\label{alg:game}
\begin{algorithmic}
\State \textbf{Initialize} $N$ policy networks $\{\mathcal{A}_{\phi_i}^i\}_{i=1}^N$ and $N$ value networks $\{V_{\xi_i}^i\}_{i=1}^N$.
\State \textbf{Initialize} hyperparameters: learning rate $\eta$, loss weights $\{\omega_i\}_{i=1}^N$.

\For{training iteration $k=1, 2, \dots$}
    \State Sample a batch of state-time points $\{(t_j, \statevec_j)\}_{j=1}^B$.
    \State Let $\mathcal{L}_{\text{Game}} \leftarrow 0$.
    
    \State \textbf{Compute Policies for All Players at all points $j$:}
    \State For $p=1,\dots,N$, compute $a_j^p \leftarrow \mathcal{A}_{\phi_p}^p(t_j, \statevec_j)$ for all $j=1,\dots,B$.
    
    \For{each player $i=1, \dots, N$}
        \State \textbf{Compute Derivatives for Player $i$:}
        \State Compute $\partial_t V_{\xi_i}^i(t_j,\statevec_j), \nabla_{\statevec} V_{\xi_i}^i(t_j,\statevec_j), \nabla_{\statevec\statevec}^2 V_{\xi_i}^i(t_j,\statevec_j)$ via AD.
        
        \State \textbf{Compute Player $i$'s Hamiltonian for each point $j$:}
        \State $H_j^i \leftarrow f^i(\statevec_j, a_j^i, a_j^{-i}) + \mathcal{D}^{a_j^i, a_j^{-i}}V_{\xi_i}^i(t_j, \statevec_j)$.
        
        \State \textbf{Compute Player $i$'s Loss Components:}
        \State $L_{\text{PERF}}^i \leftarrow -\frac{1}{B}\sum_{j=1}^B H_j^i$
        \State $L_{\text{HJB}}^i \leftarrow \frac{1}{B}\sum_{j=1}^B \left( -\partial_t V_{\xi_i}^i(t_j,\statevec_j) - H_j^i \right)^2$
        
        \State \textbf{Add to Total Loss:} $\mathcal{L}_{\text{Game}} \leftarrow \mathcal{L}_{\text{Game}} + L_{\text{PERF}}^i + \omega_i L_{\text{HJB}}^i$.
    \EndFor
    
    \State \textbf{Update All Networks Simultaneously:}
    \State Update all parameters $(\{\phi_i\}, \{\xi_i\})$ using one step of gradient descent on $\mathcal{L}_{\text{Game}}$:
    \State $(\{\phi_i\}, \{\xi_i\}) \leftarrow (\{\phi_i\}, \{\xi_i\}) - \eta \nabla_{(\{\phi_i\}, \{\xi_i\})} \mathcal{L}_{\text{Game}}$
\EndFor
\State \textbf{Return} trained networks $(\{\mathcal{A}_{\phi_i}^i\}, \{V_{\xi_i}^i\})$, which approximate the MPNE.
\end{algorithmic}
\end{algorithm}

\end{document}
